% Options for packages loaded elsewhere
\PassOptionsToPackage{unicode}{hyperref}
\PassOptionsToPackage{hyphens}{url}
%
\documentclass[
]{article}
\usepackage{amsmath,amssymb}
\usepackage{lmodern}
\usepackage{iftex}
\ifPDFTeX
  \usepackage[T1]{fontenc}
  \usepackage[utf8]{inputenc}
  \usepackage{textcomp} % provide euro and other symbols
\else % if luatex or xetex
  \usepackage{unicode-math}
  \defaultfontfeatures{Scale=MatchLowercase}
  \defaultfontfeatures[\rmfamily]{Ligatures=TeX,Scale=1}
\fi
% Use upquote if available, for straight quotes in verbatim environments
\IfFileExists{upquote.sty}{\usepackage{upquote}}{}
\IfFileExists{microtype.sty}{% use microtype if available
  \usepackage[]{microtype}
  \UseMicrotypeSet[protrusion]{basicmath} % disable protrusion for tt fonts
}{}
\makeatletter
\@ifundefined{KOMAClassName}{% if non-KOMA class
  \IfFileExists{parskip.sty}{%
    \usepackage{parskip}
  }{% else
    \setlength{\parindent}{0pt}
    \setlength{\parskip}{6pt plus 2pt minus 1pt}}
}{% if KOMA class
  \KOMAoptions{parskip=half}}
\makeatother
\usepackage{xcolor}
\usepackage{longtable,booktabs,array}
\usepackage{calc} % for calculating minipage widths
% Correct order of tables after \paragraph or \subparagraph
\usepackage{etoolbox}
\makeatletter
\patchcmd\longtable{\par}{\if@noskipsec\mbox{}\fi\par}{}{}
\makeatother
% Allow footnotes in longtable head/foot
\IfFileExists{footnotehyper.sty}{\usepackage{footnotehyper}}{\usepackage{footnote}}
\makesavenoteenv{longtable}
\usepackage{graphicx}
\makeatletter
\def\maxwidth{\ifdim\Gin@nat@width>\linewidth\linewidth\else\Gin@nat@width\fi}
\def\maxheight{\ifdim\Gin@nat@height>\textheight\textheight\else\Gin@nat@height\fi}
\makeatother
% Scale images if necessary, so that they will not overflow the page
% margins by default, and it is still possible to overwrite the defaults
% using explicit options in \includegraphics[width, height, ...]{}
\setkeys{Gin}{width=\maxwidth,height=\maxheight,keepaspectratio}
% Set default figure placement to htbp
\makeatletter
\def\fps@figure{htbp}
\makeatother
\usepackage[normalem]{ulem}
\setlength{\emergencystretch}{3em} % prevent overfull lines
\providecommand{\tightlist}{%
  \setlength{\itemsep}{0pt}\setlength{\parskip}{0pt}}
\setcounter{secnumdepth}{-\maxdimen} % remove section numbering
\ifLuaTeX
  \usepackage{selnolig}  % disable illegal ligatures
\fi
\IfFileExists{bookmark.sty}{\usepackage{bookmark}}{\usepackage{hyperref}}
\IfFileExists{xurl.sty}{\usepackage{xurl}}{} % add URL line breaks if available
\urlstyle{same} % disable monospaced font for URLs
\hypersetup{
  hidelinks,
  pdfcreator={LaTeX via pandoc}}

\author{}
\date{}

\begin{document}


\includegraphics[width=1.48in,height=1.48in]{vertopal_fe3807bd2608410aa17783721e5463e8/media/image1.jpeg}

\begin{quote}
\textbf{BỘ GIÁO DỤC VÀ ĐÀO TẠO}

\textbf{TRƯỜNG ĐẠI HỌC NÔNG LÂM TP HCM}

\textbf{KHOA CÔNG NGHỆ THÔNG TIN}
\end{quote}

TIỂU LUẬN TỐT NGHIỆP

\begin{quote}
\textbf{TÊN ĐỀ TÀI}
\end{quote}

\textbf{Nghiên cứu cải tiến TickNets bằng cơ chế}

\textbf{attention CBAM cho bài toán phân loại ảnh}

\begin{quote}
\textbf{Ngành : Công nghệ thông tin.}

\textbf{Niên khoá : 2022 -- 2026.}

\textbf{Lớp : DH22DTC.}

\textbf{Sinh viên thực hiện : Chu Toàn Đức.}
\end{quote}


\includegraphics[width=1.48in,height=1.47833in]{vertopal_fe3807bd2608410aa17783721e5463e8/media/image1.jpeg}

\begin{quote}
\textbf{BỘ GIÁO DỤC VÀ ĐÀO TẠO}

\textbf{TRƯỜNG ĐẠI HỌC NÔNG LÂM TP HCM}

\textbf{KHOA CÔNG NGHỆ THÔNG TIN}

TIỂU LUẬN TỐT NGHIỆP

\textbf{TÊN ĐỀ TÀI}
\end{quote}

\textbf{Nghiên cứu cải tiến TickNets bằng cơ chế}

\textbf{attention CBAM cho bài toán phân loại ảnh}

\begin{quote}
\textbf{CÁN BỘ HƯỚNG DẪN} \textbf{SINH VIÊN THỰC HIỆN}

TS. Nguyễn Văn Dũ 1. Chu Toàn Đức (22130047).

TP.HỒ CHÍ MINH,24 tháng 07 năm 2026
\end{quote}

\hypertarget{danh-suxe1ch-chux1eef-viux1ebft-tux1eaft}{%
\section{DANH SÁCH CHỮ VIẾT
TẮT}\label{danh-suxe1ch-chux1eef-viux1ebft-tux1eaft}}

\begin{quote}
NLP \textbf{N}atural \textbf{L}anguage \textbf{P}rocessing Xử lý ngôn
ngữ tự nhiên.

CNN \textbf{C}onvolutional \textbf{N}eural \textbf{N}etwork Mạng nơ-ron
tích chập.
\end{quote}

CV \textbf{C}omputer \textbf{V}ision.

Thị giác máy tính.

\begin{quote}
DL \textbf{D}eep \textbf{L}earning Học sâu.

ANN \textbf{A}rtificial \textbf{N}eural \textbf{N}etwork Mạng nơ-ron
nhân tạo.

FC \textbf{F}ully \textbf{C}onnected Layer Lớp đầy đủ.

CAM \textbf{C}hannel \textbf{A}ttention \textbf{M}odule Mô-đun chú ý
theo kênh.

SAM \textbf{S}patial \textbf{A}ttention \textbf{M}odule

Mô-đun chú ý theo không gian.

KNN \textbf{K}-\textbf{N}earest \textbf{N}eighbors K láng giềng gần
nhất.
\end{quote}

\hypertarget{danh-mux1ee5c-huxecnh-ux1ea3nh}{%
\section{DANH MỤC HÌNH ẢNH}\label{danh-mux1ee5c-huxecnh-ux1ea3nh}}

\begin{quote}
Hình 2.1. Minh họa về Neural Network 5

Hình 2.2 Ví dụ về kiến trúc của CNN 6

Hình 2.3 Ví dụ về Output của Convolution 8

Hình 2.4 Ví dụ về việc sử dụng hàm kích hoạt ReLU 9

Hình 2.4 Ví dụ về Max-pooling và Average-pooling 11

Hình 2.5 Minh họa Fully Connected Layer 12

Hình 2.6 Kiến trúc mô hình VGG16 13

Hình 2.7 Ví dụ cách thức hoạt động của CNN 14

Hình 2.8 Công thức tính toán output của Convolution 15

Hình 2.9 Mô hình hoạt động của CAM 17

Hình 2.10 Mô hình hoạt động của SAM 19

Hình 2.11 Mô hình CBAM 20

Hình 3.1 Ví dụ về bài toán phân loại ảnh 22

Hình 3.2 Mô hình CNN kết hợp CBAM 23

Hình 3.3 Quá trình xây dựng và huấn luyện mô hình 24

Hình 3.4 Tải tập dữ liệu 26

Hình 3.5 Mô hình CNN cho tập dữ liệu MNIST-Fashion 28

Hình 3.6 Mô hình CNN cho tập dữ liệu Cifar-10 29

Hình 3.7 Cấu trúc Confusion matrix 31

Hình 3.8 Biểu đồ loss và accuracy của mô hình CNN có áp dụng Attention
và không áp dụng với tập dữ liệu MNIST-Fashion 33

Hình 3.9 Biểu đồ loss và accuracy của mô hình CNN có áp dụng Attention
và không áp dụng với tập dữ liệu Cifar-10 34
\end{quote}

\hypertarget{danh-mux1ee5c-bux1ea3ng}{%
\section{DANH MỤC BẢNG}\label{danh-mux1ee5c-bux1ea3ng}}

\begin{quote}
Bảng 3.1 Bảng mô tả dữ liệu tập dữ liệu MNIST-Fashion 25

Bảng 3.2 Bảng mô tả dữ liệu tập dữ liệu Cifar-10 26

Bảng 3.3 Bảng so sánh kết quả giữa các giải thuật trên tập dữ liệu
MNIST-Fashion

. 35

Bảng 3.4 Bảng so sánh kết quả giữa các giải thuật trên tập dữ liệu
Cifar-10 35

\textbf{TÓM TẮT}

Bài báo cáo này nghiên cứu, đề xuất và hiện thực hóa giải pháp cải tiến
kiến trúc mạng nơ-ron tích chập nhẹ TickNets bằng cách tích hợp phân cấp
cơ chế chú ý khối tích chập CBAM (Convolutional Block Attention Module)
cho bài toán phân loại hình ảnh. Thông qua các thử nghiệm thực tiễn trên
hai nhóm tập dữ liệu chuẩn học thuật Fashion-MNIST, CIFAR-10 và ứng dụng
thực tế phức tạp PlantVillage, Chest-Xray-Pneumonia, kết quả thực nghiệm
chỉ ra rằng giải pháp chèn CBAM phân cấp tại các điểm nối mạng không chỉ
giúp loại bỏ hiệu quả nhiễu nền trên ảnh bối cảnh thực tế mà còn khắc
phục tối ưu sự suy giảm hiệu năng cục bộ trên ảnh kích thước nhỏ. Dựa
trên các độ đo chuẩn mực bao gồm accuracy, precision, recall và
f1-score, mô hình cải tiến đề xuất mang lại hiệu suất vượt trội và ổn
định hơn đáng kể so với các mô hình baseline không áp dụng hoặc chèn cơ
chế chú ý sai vị trí, đóng góp một hướng đi có giá trị khoa học cao cho
các ứng dụng học sâu tối giản trên thiết bị có tài nguyên hạn chế.
\end{quote}

\hypertarget{mux1ee5c-lux1ee5c}{%
\section{MỤC LỤC}\label{mux1ee5c-lux1ee5c}}

\protect\hyperlink{danh-suxe1ch-chux1eef-viux1ebft-tux1eaft}{DANH SÁCH
CHỮ VIẾT TẮT
\protect\hyperlink{danh-suxe1ch-chux1eef-viux1ebft-tux1eaft}{1}}

\protect\hyperlink{danh-mux1ee5c-huxecnh-ux1ea3nh}{DANH MỤC HÌNH ẢNH
\protect\hyperlink{danh-mux1ee5c-huxecnh-ux1ea3nh}{2}}

\protect\hyperlink{danh-mux1ee5c-bux1ea3ng}{DANH MỤC BẢNG
\protect\hyperlink{danh-mux1ee5c-bux1ea3ng}{3}}

\protect\hyperlink{_Toc235870340}{MỤC LỤC
\protect\hyperlink{_Toc235870340}{5}}

\protect\hyperlink{mux1edf-ux111ux1ea7u}{MỞ ĐẦU
\protect\hyperlink{mux1edf-ux111ux1ea7u}{7}}

\protect\hyperlink{luxfd-do-chux1ecdn-ux111ux1ec1-tuxe0i}{1. LÝ DO CHỌN
ĐỀ TÀI: \protect\hyperlink{luxfd-do-chux1ecdn-ux111ux1ec1-tuxe0i}{7}}

\protect\hyperlink{mux1ee5c-tiuxeau-vuxe0-phux1ea1m-vi-nghiuxean-cux1ee9u}{2.
MỤC TIÊU VÀ PHẠM VI NGHIÊN CỨU
\protect\hyperlink{mux1ee5c-tiuxeau-vuxe0-phux1ea1m-vi-nghiuxean-cux1ee9u}{7}}

\protect\hyperlink{uxfd-nghux129a-khoa-hux1ecdc-vuxe0-thux1ef1c-tiux1ec5n-cux1ee7a-ux111ux1ec1-tuxe0i}{3.
Ý NGHĨA KHOA HỌC VÀ THỰC TIỄN CỦA ĐỀ TÀI
\protect\hyperlink{uxfd-nghux129a-khoa-hux1ecdc-vuxe0-thux1ef1c-tiux1ec5n-cux1ee7a-ux111ux1ec1-tuxe0i}{8}}

\protect\hyperlink{chux1b0ux1a1ng-1-tux1ed5ng-quan-ux111ux1ec1-tuxe0i}{CHƯƠNG
1 TỔNG QUAN ĐỀ TÀI
\protect\hyperlink{chux1b0ux1a1ng-1-tux1ed5ng-quan-ux111ux1ec1-tuxe0i}{10}}

\protect\hyperlink{chux1b0ux1a1ng-2-cux1a1-sux1edf-luxfd-thuyux1ebft}{CHƯƠNG
2 CƠ SỞ LÝ THUYẾT
\protect\hyperlink{chux1b0ux1a1ng-2-cux1a1-sux1edf-luxfd-thuyux1ebft}{11}}

\protect\hyperlink{mux1ea1ng-nux1a1-ron-tuxedch-chux1eadp-nhux1eb9-siuxeau-nhux1eb9-ticknets}{2.1.
MẠNG NƠ-RON TÍCH CHẬP NHẸ SIÊU NHẸ TICKNETS:
\protect\hyperlink{mux1ea1ng-nux1a1-ron-tuxedch-chux1eadp-nhux1eb9-siuxeau-nhux1eb9-ticknets}{11}}

\protect\hyperlink{mux1ea1ng-tuxedch-chux1eadp-siuxeau-nhux1eb9-ticknets}{2.1.1.
Mạng tích chập siêu nhẹ TickNets
\protect\hyperlink{mux1ea1ng-tuxedch-chux1eadp-siuxeau-nhux1eb9-ticknets}{11}}

\protect\hyperlink{cuxe1c-khuxe1i-niux1ec7m-tuxedch-chux1eadp-tux1ed1i-giux1ea3n.}{2.1.2
Các khái niệm tích chập tối giản.
\protect\hyperlink{cuxe1c-khuxe1i-niux1ec7m-tuxedch-chux1eadp-tux1ed1i-giux1ea3n.}{12}}

\protect\hyperlink{kiux1ebfn-truxfac-khux1ed1i-full-residual-point-depth-point}{2.1.3
Kiến trúc khối Full-Residual Point-Depth-Point
\protect\hyperlink{kiux1ebfn-truxfac-khux1ed1i-full-residual-point-depth-point}{13}}

\protect\hyperlink{cux1ea5u-truxfac-xux1b0ux1a1ng-sux1ed1ng-co-giuxe3n-kuxeanh-huxecnh-dux1ea5u-tuxedch-vuxe0-muxf4-huxecnh-ticknet-small}{2.1.4
Cấu trúc xương sống co giãn kênh hình dấu tích và mô hình TickNet-small
\protect\hyperlink{cux1ea5u-truxfac-xux1b0ux1a1ng-sux1ed1ng-co-giuxe3n-kuxeanh-huxecnh-dux1ea5u-tuxedch-vuxe0-muxf4-huxecnh-ticknet-small}{15}}

\protect\hyperlink{cux1a1-chux1ebf-lan-truyux1ec1n-ux111ux1eb7c-trux1b0ng-vuxe0-cuxe1ch-thux1ee9c-houx1ea1t-ux111ux1ed9ng-cux1ee7a-ticknets}{2.1.5
Cơ chế lan truyền đặc trưng và cách thức hoạt động của TickNets
\protect\hyperlink{cux1a1-chux1ebf-lan-truyux1ec1n-ux111ux1eb7c-trux1b0ng-vuxe0-cuxe1ch-thux1ee9c-houx1ea1t-ux111ux1ed9ng-cux1ee7a-ticknets}{16}}

\protect\hyperlink{cux1a1-chux1ebf-chuxfa-uxfd-attention-mechanism}{2.2
CƠ CHẾ CHÚ Ý (ATTENTION MECHANISM):
\protect\hyperlink{cux1a1-chux1ebf-chuxfa-uxfd-attention-mechanism}{17}}

\protect\hyperlink{squeeze-and-excitation-mux1eb7c-ux111ux1ecbnh-trong-ticknets}{2.2.1
Squeeze-and-Excitation mặc định trong TickNets
\protect\hyperlink{squeeze-and-excitation-mux1eb7c-ux111ux1ecbnh-trong-ticknets}{17}}

\protect\hyperlink{channel-attention-module-cam}{2.2.2 Channel Attention
Module (CAM): \protect\hyperlink{channel-attention-module-cam}{19}}

\protect\hyperlink{spatial-attention-module-sam}{2.2.3 Spatial Attention
Module (SAM): \protect\hyperlink{spatial-attention-module-sam}{20}}

\protect\hyperlink{muxf4-ux111un-chuxfa-uxfd-khux1ed1i-tuxedch-chux1eadp-cbam-houx1ea1t-ux111ux1ed9ng-tuux1ea7n-tux1ef1-cam-sam}{2.2.4
Mô-đun chú ý khối tích chập CBAM hoạt động tuần tự CAM + SAM
\protect\hyperlink{muxf4-ux111un-chuxfa-uxfd-khux1ed1i-tuxedch-chux1eadp-cbam-houx1ea1t-ux111ux1ed9ng-tuux1ea7n-tux1ef1-cam-sam}{21}}

\protect\hyperlink{phux1b0ux1a1ng-phuxe1p-nghiuxean-cux1ee9u}{2.3.
PHƯƠNG PHÁP NGHIÊN CỨU
\protect\hyperlink{phux1b0ux1a1ng-phuxe1p-nghiuxean-cux1ee9u}{23}}

\protect\hyperlink{chux1b0ux1a1ng-3-buxe0i-touxe1n-phuxe2n-loux1ea1i-huxecnh-ux1ea3nh-vuxe0-ux111ux1ec1-xuux1ea5t-cux1ea3i-tiux1ebfn}{CHƯƠNG
3 BÀI TOÁN PHÂN LOẠI HÌNH ẢNH VÀ ĐỀ XUẤT CẢI TIẾN
\protect\hyperlink{chux1b0ux1a1ng-3-buxe0i-touxe1n-phuxe2n-loux1ea1i-huxecnh-ux1ea3nh-vuxe0-ux111ux1ec1-xuux1ea5t-cux1ea3i-tiux1ebfn}{25}}

\protect\hyperlink{giux1edbi-thiux1ec7u-buxe0i-touxe1n-phuxe2n-loux1ea1i-huxecnh-ux1ea3nh}{3.1
GIỚI THIỆU BÀI TOÁN PHÂN LOẠI HÌNH ẢNH
\protect\hyperlink{giux1edbi-thiux1ec7u-buxe0i-touxe1n-phuxe2n-loux1ea1i-huxecnh-ux1ea3nh}{25}}

\protect\hyperlink{ux111ux1ec1-xuux1ea5t-muxf4-huxecnh-ticknets-kux1ebft-hux1ee3p-cux1a1-chux1ebf-chuxfa-uxfd-phuxe2n-cux1ea5p-hierarchical-attention}{3.2
ĐỀ XUẤT MÔ HÌNH TICKNETS KẾT HỢP CƠ CHẾ CHÚ Ý PHÂN CẤP HIERARCHICAL
ATTENTION
\protect\hyperlink{ux111ux1ec1-xuux1ea5t-muxf4-huxecnh-ticknets-kux1ebft-hux1ee3p-cux1a1-chux1ebf-chuxfa-uxfd-phuxe2n-cux1ea5p-hierarchical-attention}{25}}

\protect\hyperlink{thiux1ebft-kux1ebf-kiux1ebfn-truxfac-ux111ux1ec1-xuux1ea5t-ticknet-small-cbam-hook}{3.2.1.
Thiết kế kiến trúc đề xuất TickNet-small CBAM-Hook
\protect\hyperlink{thiux1ebft-kux1ebf-kiux1ebfn-truxfac-ux111ux1ec1-xuux1ea5t-ticknet-small-cbam-hook}{26}}

\protect\hyperlink{cux1a1-chux1ebf-chuxfa-uxfd-phuxe2n-cux1ea5p}{3.2.2
Cơ chế chú ý phân cấp
\protect\hyperlink{cux1a1-chux1ebf-chuxfa-uxfd-phuxe2n-cux1ea5p}{26}}

\protect\hyperlink{phuxe2n-tuxedch-touxe1n-hux1ecdc-giux1ea3i-quyux1ebft-sux1ef1-chux1ed3ng-chuxe9o-ux111ux1eb7c-trux1b0ng-khuxf4ng-gian-giux1eefa-depthwise-convolution-vuxe0-spatial-attention-module}{3.2.3
Phân tích toán học giải quyết sự chồng chéo đặc trưng không gian giữa
Depthwise Convolution và Spatial Attention Module
\protect\hyperlink{phuxe2n-tuxedch-touxe1n-hux1ecdc-giux1ea3i-quyux1ebft-sux1ef1-chux1ed3ng-chuxe9o-ux111ux1eb7c-trux1b0ng-khuxf4ng-gian-giux1eefa-depthwise-convolution-vuxe0-spatial-attention-module}{27}}

\protect\hyperlink{hiux1ec7n-thux1ef1c-giux1ea3i-phuxe1p}{3.3 HIỆN THỰC
GIẢI PHÁP:
\protect\hyperlink{hiux1ec7n-thux1ef1c-giux1ea3i-phuxe1p}{28}}

\protect\hyperlink{muxf4i-trux1b0ux1eddng-thux1ef1c-nghiux1ec7m-vuxe0-xuxe2y-dux1ef1ng-muxf4-huxecnh}{3.3.1
Môi trường thực nghiệm và xây dựng mô hình:
\protect\hyperlink{muxf4i-trux1b0ux1eddng-thux1ef1c-nghiux1ec7m-vuxe0-xuxe2y-dux1ef1ng-muxf4-huxecnh}{28}}

\protect\hyperlink{kux1ebft-quux1ea3-thux1ef1c-nghiux1ec7m-vuxe0-ux111uxe1nh-giuxe1}{3.3.2
Kết quả thực nghiệm và đánh giá:
\protect\hyperlink{kux1ebft-quux1ea3-thux1ef1c-nghiux1ec7m-vuxe0-ux111uxe1nh-giuxe1}{28}}

\protect\hyperlink{chux1b0ux1a1ng-4-kux1ebft-quux1ea3-kux1ebft-luux1eadn-vuxe0-kiux1ebfn-nghux1ecb}{CHƯƠNG
4: KẾT QUẢ, KẾT LUẬN VÀ KIẾN NGHỊ
\protect\hyperlink{chux1b0ux1a1ng-4-kux1ebft-quux1ea3-kux1ebft-luux1eadn-vuxe0-kiux1ebfn-nghux1ecb}{29}}

\protect\hyperlink{kux1ebft-quux1ea3}{4.1 KẾT QUẢ:
\protect\hyperlink{kux1ebft-quux1ea3}{29}}

\protect\hyperlink{kux1ebft-luux1eadn}{4.2 KẾT LUẬN
\protect\hyperlink{kux1ebft-luux1eadn}{29}}

\protect\hyperlink{kiux1ebfn-nghux1ecb}{4.3 KIẾN NGHỊ
\protect\hyperlink{kiux1ebfn-nghux1ecb}{29}}

\protect\hyperlink{tuxe0i-liux1ec7u-tham-khux1ea3o}{TÀI LIỆU THAM KHẢO
\protect\hyperlink{tuxe0i-liux1ec7u-tham-khux1ea3o}{31}}

\hypertarget{mux1edf-ux111ux1ea7u}{%
\section{MỞ ĐẦU}\label{mux1edf-ux111ux1ea7u}}

\hypertarget{luxfd-do-chux1ecdn-ux111ux1ec1-tuxe0i}{%
\subsection{LÝ DO CHỌN ĐỀ
TÀI:}\label{luxfd-do-chux1ecdn-ux111ux1ec1-tuxe0i}}

\begin{quote}
Phân loại hình ảnh (image classification) là một trong những bài toán cơ
bản trong thị giác máy tính. Bài toán này được phát biểu như sau: cho
một hình ảnh, cần xác định hình ảnh này thuộc lớp nào trong một số lớp
cho trước. Ngày nay nó được áp dụng rộng rãi trong hầu hết các lĩnh vực
như quản lý chất lượng sản phẩm, ô tô tự lái, y tế, \ldots vì vậy việc
nâng cao tính chính xác của việc phân loại hình ảnh là một trong những
vấn đề cấp thiết.

Trong TickNet là một kiến trúc CNN nhẹ được thiết kế nhằm giảm số lượng
tham số và chi phí tính toán nhưng vẫn duy trì hiệu quả phân loại cao
thông qua khối Full-Residual Point-Depth-Point (FR-PDP) và cấu trúc
Tick-shape Backbone. Tuy nhiên, mô hình TickNet gốc chưa khai thác cơ
chế Attention để tăng cường khả năng tập trung vào các đặc trưng quan
trọng của ảnh.

Trong số các cơ chế Attention, Convolutional Block Attention Module
(CBAM) là một mô-đun nhẹ, có khả năng cải thiện việc trích xuất đặc
trưng bằng cách kết hợp chú ý theo kênh và chú ý theo không gian. Vì
vậy, việc nghiên cứu tích hợp CBAM vào TickNet-Small nhằm nâng cao hiệu
quả phân loại ảnh là một hướng nghiên cứu có ý nghĩa cả về mặt học thuật
và ứng dụng thực tiễn.

Xuất phát từ những lý do trên, đề tài "Nghiên cứu cải tiến TickNet-Small
bằng cơ chế Attention CBAM cho bài toán phân loại ảnh" được lựa chọn
thực hiện.
\end{quote}

\hypertarget{mux1ee5c-tiuxeau-vuxe0-phux1ea1m-vi-nghiuxean-cux1ee9u}{%
\subsection{MỤC TIÊU VÀ PHẠM VI NGHIÊN
CỨU}\label{mux1ee5c-tiuxeau-vuxe0-phux1ea1m-vi-nghiuxean-cux1ee9u}}

\begin{quote}
Mục tiêu của đề tài hướng đến việc nghiên cứu kiến trúc TickNet-Small và
cơ chế Attention CBAM, từ đó đề xuất mô hình TickNet-Small kết hợp CBAM
nhằm nâng cao hiệu quả của bài toán phân loại ảnh. Đồng thời, tiến hành
thực nghiệm trên nhiều bộ dữ liệu khác nhau để đánh giá khả năng cải
thiện hiệu suất của mô hình so với TickNet-Small gốc.

Phạm vi nghiên cứu:

Nghiên cứu kiến trúc TickNet, khối FR-PDP và cấu trúc Tick-shape
Backbone.

Nghiên cứu cơ chế Attention, đặc biệt là Convolutional Block Attention
Module (CBAM).

Xây dựng mô hình TickNet-Small kết hợp CBAM.

Huấn luyện và đánh giá mô hình trên các tập dữ liệu MNIST, CIFAR-10,
PlantVillage và ChestXray.

Quá trình huấn luyện sử dụng bộ tối ưu hóa SGD kết hợp với cấu hình hàm
gọi lại động bao gồm giảm tốc độ học ReduceLROnPlateau và dừng huấn
luyện sớm EarlyStopping.

Đánh giá hiệu năng dựa trên các độ đo tiêu chuẩn: Accuracy, Precision,
Recall, F1-score cùng trực quan hóa bằng Confusion Matrix và bản đồ
nhiệt Grad-CAM để giải thích vùng chú ý của mô hình.
\end{quote}

\hypertarget{uxfd-nghux129a-khoa-hux1ecdc-vuxe0-thux1ef1c-tiux1ec5n-cux1ee7a-ux111ux1ec1-tuxe0i}{%
\subsection{Ý NGHĨA KHOA HỌC VÀ THỰC TIỄN CỦA ĐỀ
TÀI}\label{uxfd-nghux129a-khoa-hux1ecdc-vuxe0-thux1ef1c-tiux1ec5n-cux1ee7a-ux111ux1ec1-tuxe0i}}

\begin{quote}
Ý nghĩa khoa học:
\end{quote}

\begin{itemize}
\item
  Giải quyết xung đột kiến trúc trong mạng siêu nhẹ: Đề tài đóng góp một
  lời giải khoa học sâu sắc cho việc tối ưu hóa hiệu năng các mạng
  nơ-ron siêu nhẹ bằng cách chỉ ra và xử lý thành công sự chồng chéo
  thông tin giữa phép tích chập chiều sâu và mô-đun chú ý không gian
  trong các cấu trúc Perceptron cục bộ.
\item
  Đề xuất mô hình chú ý phân cấp mới: Nghiên cứu tiên phong trong việc
  thiết lập cấu trúc phối hợp hai tầng cơ chế chú ý: giữ nguyên
  attention kênh SE ở cấp độ khối cơ bản để lọc đặc trưng cục bộ, và áp
  dụng attention hỗn hợp CBAM ở cấp độ điểm nối để kiểm soát đặc trưng
  toàn cục. Điều này mở ra một hướng tiếp cận mới mẻ trong việc thiết kế
  và tích hợp các khối attention vào các mạng nơ-ron tối giản.
\item
  Góp phần phát triển lý thuyết học sâu giải thích được : Thông qua việc
  ứng dụng thuật toán Grad-CAM để trực quan hóa các bản đồ nhiệt chú ý ,
  đề tài cung cấp các minh chứng khoa học trực quan, rõ ràng giải thích
  tại sao mô hình cải tiến có thể hội tụ nhanh hơn, định vị chính xác
  hơn các vùng bệnh lý hoặc vật thể so với mô hình baseline thông
  thường\ldots{}
\end{itemize}

\begin{quote}
Ý nghĩa thực tiễn: kết quả nghiên cứu trình bày trong tiểu luận có ý
nghĩa rất lớn khi áp dụng vào các bài toán trong thực tiễn như sau
\end{quote}

\begin{itemize}
\item
  Nâng cao chất lượng chẩn đoán y khoa và quản lý nông nghiệp: Kết quả
  thực nghiệm của mô hình cải tiến trên tập dữ liệu Chest-Xray-Pneumonia
  giúp tối ưu hóa độ chính xác và giảm thiểu tỷ lệ bỏ sót trong việc
  chẩn đoán viêm phổi qua hình ảnh X-quang phổi . Đồng thời, thực nghiệm
  trên tập PlantVillage giúp hỗ trợ bà con nông dân nhận diện sớm và
  chính xác các loại bệnh trên lá cây trồng, góp phần thúc đẩy các giải
  pháp nông nghiệp thông minh, công nghệ cao.
\item
  Tính khả thi cao trong triển khai thực tế trên thiết bị biên: Nhờ kế
  thừa ưu điểm siêu nhẹ của mạng TickNets, mô hình TickNet-small cải
  tiến bằng CBAM phân cấp cực kỳ phù hợp để đóng gói và vận hành trực
  tiếp trên các thiết bị di động cá nhân, máy tính bảng cầm tay, thiết
  bị bay không người lái hay các chip nhúng biên IoT có năng lượng và
  năng lực tính toán hạn chế mà không cần phụ thuộc vào máy chủ đám mây.
\item
\end{itemize}

\hypertarget{chux1b0ux1a1ng-1-tux1ed5ng-quan-ux111ux1ec1-tuxe0i}{%
\section{CHƯƠNG 1 TỔNG QUAN ĐỀ
TÀI}\label{chux1b0ux1a1ng-1-tux1ed5ng-quan-ux111ux1ec1-tuxe0i}}

\begin{quote}
Cơ chế chú ý (Attention Mechanism) là một trong những cơ chế được sử
dụng phổ biến trong các bài toán xử lý ngôn ngữ tự nhiên (NLP), đặc biệt
là trong các mô hình dịch máy. Một trong những thành tựu nổi bật của nó
là sự phát triển của mô hình Transformer, giới thiệu trong bài báo
``Attention is All You Need'' của Vaswani et al. năm 2017 {[}1{]}. Mô
hình Transformer đã thay đổi cách tiếp cận dịch máy và các bài toán xử
lý ngôn ngữ tự nhiên (NLP) nhờ vào cơ chế Attention mạnh mẽ của nó. Góp
phần dẫn đến sự phát triển của các mô hình như BERT (Bidirectional
Encoder Representations from Transformers) và GPT (Generative
Pre-trained Transformer). Những mô hình này đã thiết lập các chuẩn mực
mới trong nhiều bài toán NLP khác nhau, bao gồm dịch máy, hiểu ngôn ngữ
tự nhiên, và tạo văn bản.

Mặc dù các nghiên cứu tích hợp cơ chế chú ý như CBAM đã chứng minh hiệu
quả vượt trội khi kết hợp với các mạng tích chập truyền thống như VGG16
, việc áp dụng cơ chế này lên các kiến trúc mạng siêu nhẹ vẫn đối mặt
với nhiều thách thức lớn về sự tương thích cấu trúc . Các mạng siêu nhẹ
phổ biến hiện nay như MobileNets hay ShuffleNets thường chỉ sử dụng cơ
chế chú ý một hướng SE để tái cấu trúc kênh đặc trưng nhằm tiết kiệm
tham số . Kiến trúc mạng siêu nhẹ cải tiến gần đây là TickNets -- nổi
bật với khối FR-PDP và cấu trúc xương sống co giãn kênh hình dấu tích
tuy đạt hiệu suất cao nhưng cũng mới chỉ dừng lại ở việc tích hợp mặc
định cơ chế attention SE ở mức cục bộ bên trong khối.

Do đó, bài nghiên cứu này tập trung vào giải pháp cải tiến mạng
TickNet-small bằng cơ chế chú ý phân cấp. Phương pháp này đề xuất giữ
nguyên SE cục bộ trong khối FR-PDP, đồng thời chèn CBAM một cách chiến
lược tại các điểm nối mạng giữa hai backbone co giãn và trước lớp Global
Average Pooling cuối mạng , giúp mô hình phát huy tối đa sức mạnh lọc
nhiễu không gian toàn cục mà không làm bùng nổ tham số tính toán.
\end{quote}

\hypertarget{chux1b0ux1a1ng-2-cux1a1-sux1edf-luxfd-thuyux1ebft}{%
\section{\texorpdfstring{CHƯƠNG 2\\
CƠ SỞ LÝ
THUYẾT}{CHƯƠNG 2 CƠ SỞ LÝ THUYẾT}}\label{chux1b0ux1a1ng-2-cux1a1-sux1edf-luxfd-thuyux1ebft}}

\begin{enumerate}
\def\labelenumi{\arabic{enumi}.}
\item ~
  \hypertarget{mux1ea1ng-nux1a1-ron-tuxedch-chux1eadp-nhux1eb9-siuxeau-nhux1eb9-ticknets}{%
  \subsection{MẠNG NƠ-RON TÍCH CHẬP NHẸ SIÊU NHẸ
  TICKNETS:}\label{mux1ea1ng-nux1a1-ron-tuxedch-chux1eadp-nhux1eb9-siuxeau-nhux1eb9-ticknets}}

  \begin{enumerate}
  \def\labelenumii{\arabic{enumii}.}
  \item ~
    \hypertarget{mux1ea1ng-tuxedch-chux1eadp-siuxeau-nhux1eb9-ticknets}{%
    \subsubsection{Mạng tích chập siêu nhẹ
    TickNets}\label{mux1ea1ng-tuxedch-chux1eadp-siuxeau-nhux1eb9-ticknets}}
  \end{enumerate}
\end{enumerate}

\begin{quote}
Trong xu hướng phát triển của học sâu ứng dụng trên các thiết bị di động
và hệ thống nhúng, các mạng nơ-ron tích chập siêu nhẹ đóng vai trò tối
quan trọng nhờ khả năng cân bằng giữa kích thước mô hình tối giản, độ
phức tạp tính toán thấp và hiệu suất năng lượng tối ưu . Các kiến trúc
kinh điển như MobileNets, ShuffleNets hay EfficientNets đã gặt hái được
nhiều thành công nhờ thay thế các phép tích chập tiêu chuẩn bằng các
phép tích chập tối giản có chi phí thấp.

Tuy nhiên các kiến trúc siêu nhẹ truyền thống này vẫn tồn tại hai hạn
chế cấu trúc cốt lõi:

Sự bùng nổ tham số ở các tầng cuối: Backbone của các mạng này được thiết
kế theo nguyên tắc tăng dần đều số lượng kênh đặc trưng. Khi mạng càng
đi sâu, số lượng kênh ở các tầng cuối cùng tăng lên rất lớn, dẫn đến
kích thước mô hình bị phình to ra nhanh chóng.

Sự thiếu hụt liên kết đồng nhất. Dựa trên cơ chế kết nối tàn dư nhằm
giải quyết vấn đề triệt tiêu gradient , các mạng siêu nhẹ thường sử dụng
các liên kết tắt . Tuy nhiên, các liên kết tắt này chỉ là bán tàn dư ,
tức là chúng chỉ được kích hoạt khi kích thước không gian và số lượng
kênh đặc trưng giữa đầu vào và đầu ra hoàn toàn trùng khớp. Khi xảy ra
sự thay đổi kích thước không gian do bước trượt tích chập s
\textgreater{} 1 hoặc sự thay đổi số lượng kênh đặc trưng giữa các khối
nơ-ron, liên kết đồng nhất bị phá vỡ.

Để khắc phục triệt để các hạn chế trên, họ mạng TickNets được đề xuất
như một kiến trúc mạng tích chập siêu nhẹ thế hệ mới dựa trên ba đóng
góp khoa học chính: khối tích chập tối giản FR-PDP sử dụng liên kết tàn
dư đầy đủ để duy trì các liên kết đồng nhất trực tiếp trong mọi điều
kiện; cấu trúc xương sống co giãn kênh hình dấu tích nhằm tối ưu hóa
phân bổ tham số để linh hoạt tùy biến độ sâu của mô hình.
\end{quote}

\hypertarget{cuxe1c-khuxe1i-niux1ec7m-tuxedch-chux1eadp-tux1ed1i-giux1ea3n.}{%
\subsubsection{Các khái niệm tích chập tối
giản.}\label{cuxe1c-khuxe1i-niux1ec7m-tuxedch-chux1eadp-tux1ed1i-giux1ea3n.}}

\begin{quote}
Để giảm chi phí tính toán cực kỳ đắt đỏ của các phép tích chập tiêu
chuẩn vốn đòi hỏi dung lượng bộ nhớ lớn , họ mạng TickNets tận dụng các
phép tích chập tối giản, bao gồm tích chập chiều sâu và tích chập điểm.

\textbf{Phép tích chập tiêu chuẩn} \textbf{:} Cho một tensor đầu vào
\(\mathcal{T \in}\mathbb{R}^{C_{in} \times H \times W}\) với \(C_{in}\)
là số lượng kênh đầu vào, H\(\times\)W là kích thước không gian của mỗi
kênh đặc trưng. Phép tiêu chuẩn biến đổi \(\mathcal{T}\) thành một
tensor đầu ra
\(\mathcal{T}_{\mu} \in \mathbb{R}^{C_{out} \times H \times W}\) bằng
cách áp dụng \(C_{out}\) bộ lọc kernels kích thước
\(C_{in} \times K\  \times K\ :\)
\end{quote}

\[\mathcal{T}_{\mu}\  = \mu_{i}\mathcal{*T,\ }i\ \  = \ 1,\ \ldots,\ C_{out}\]

\begin{quote}
Chi phí tính toán của phép toán này là cực kỳ lớn và tỷ lệ thuận với
tích của cả số kênh đầu vào và đầu ra:
\end{quote}

\[\mathcal{O}\left( H\  \times W\  \times C_{in} \times C_{out} \times K\  \times K \right)\]

\begin{quote}
\textbf{Phép tích chập chiều sâu :} Để tối giản hóa quá trình trích xuất
đặc trưng không gian, phép tích chập chiều sâu tách biệt quá trình lọc
theo từng kênh đặc trưng \(c_{j}\) của tensor đầu vào \(\mathcal{T}\)
bằng cách sử dụng các bộ lọc độc lập kích thước
\(1 \times K \times K\ \):
\end{quote}

\[\mathcal{T}_{d} = \xi_{i}*c_{j}\ ,\ \ j = 1,\ldots,C_{in}\]

\begin{quote}
Trong đó,
\(\mathcal{T}_{d} \in \mathbb{R}^{C_{in} \times H \times W}\)và
\(\xi_{i}\) đại diện cho bộ lọc tương ứng với kênh đặc trưng thứ j .
Phép toán này chỉ thực hiện vai trò lọc không gian mà không tạo ra bất
kỳ sự tương tác kênh mới nào.

\textbf{Phép tích chập điểm :} Nhằm tổng hợp và thiết lập các mối quan
hệ phi tuyến giữa các kênh đặc trưng sau khi đã được lọc không gian bởi
phép tích chập chiều sâu, phép tích chập điểm được áp dụng bằng cách sử
dụng các bộ lọc kích thước \(C_{in} \times 1 \times 1\) quét qua các
điểm ảnh :
\end{quote}

\[\mathcal{T}_{p} = \ \zeta_{j}*\ \mathcal{T}_{d},\ \ j\  = \ 1,\ \ldots,\ C_{out}\]

\begin{quote}
Hệ quả là thu được tensor đặc trưng mới
\(\mathcal{T}_{p} \in \mathbb{R}^{C_{out} \times H \times W}\). . Phép
toán này đóng vai trò quan trọng trong việc co giãn linh hoạt số lượng
kênh đặc trưng.dựa trên mục tiêu thiết kế kiến trúc.
\end{quote}

\hypertarget{kiux1ebfn-truxfac-khux1ed1i-full-residual-point-depth-point}{%
\subsubsection{Kiến trúc khối Full-Residual
Point-Depth-Point}\label{kiux1ebfn-truxfac-khux1ed1i-full-residual-point-depth-point}}

\begin{quote}
Trọng tâm cải tiến của họ mạng TickNets nằm ở thiết kế khối FR-PDP. Khối
này kết hợp mô-đun trích xuất đặc trưng phân tầng PDP độc đáo với cơ chế
liên kết tàn dư đầy đủ để khắc phục tình trạng gián đoạn liên kết đồng
nhất khi mạng thực hiện giảm chiều không gian hoặc co giãn kênh .

\textbf{Mô-đun trích xuất đặc trưng Point-Depth-Point}. Khác với khối
Bottleneck ngược của MobileNetV2 mở rộng số kênh ở đầu khối, khối PDP
thực hiện nén và giữ nguyên số lượng kênh đặc trưng ở giai đoạn đầu để
bảo toàn các thông tin đặc trưng nguyên bản. Quy trình biến đổi một
tensor đầu vào
\(\mathcal{T \in}\mathbb{R}^{C_{in} \times H \times W}\ \)qua mô đun PDP
bao gồm 4 bước liên tiếp:

Bước 1 - Tích chập điểm thứ nhất : Áp dụng \(C_{in}\) bộ lọc kích thước
\(1 \times 1\) với bước trượt \(s = 1\ \)để tạo ra các đặc trưng điểm
ảnh được nén nhưng giữ nguyên chiều sâu kênh gốc nhằm cung cấp đầu vào
chất lượng cho tích chập chiều sâu :
\end{quote}

\[\mathcal{P}_{w1} = \ \zeta_{j}\mathcal{*\ T,\ \ \ }j\  = \ 1,\ \ldots,\ C_{in}\]

\begin{quote}
Trong đó, \(\mathcal{P}_{w1} \in \mathbb{R}^{C_{in} \times H \times W}\)

Bước 2 - Tích chập chiều sâu : Sử dụng bộ lọc kích thước \(3 \times 3\)
với bước trượt \(s = 1\ \)hoặc \(s > 1\) trong trường hợp cần giảm kích
thước không gian để học các đặc trưng trong cục bộ.
\end{quote}

\[\mathcal{D}_{ws} = \xi_{j}^{3 \times 3}*c_{j} \in \mathcal{P}_{w1},\quad j = 1,\ldots,C_{in}\]

\begin{quote}
Trong đó,
\(\mathcal{D}_{ws} \in \mathbb{R}^{C_{in} \times H_{s} \times W_{s}}\).

Bước 3 - Tích chập điểm thứ 2 : Thiết lập sự tương tác kênh và co giãn
số lượng kênh đặc trưng đầu ra thành \(C_{out}\) kênh theo cấu hình
thiết kế bằng phép tích chập điểm tuần tự:
\end{quote}

\[\mathcal{P}_{w2} = \mu_{j}*\mathcal{D}_{ws},\quad j = 1,\ldots,C_{out}\]

\begin{quote}
Trong đó,
\(\mathcal{P}_{w2} \in \mathbb{R}^{C_{out} \times H_{s} \times W_{s}}\).

Bước 4 - Cơ chế chú ý kênh : Để tối ưu hóa biểu diễn đặc trưng, khối PDP
áp dụng cơ chế chú ý SE với tỷ lệ giảm mặc định \(r = 16\). Khác biệt
cốt lõi của TickNets là tính toán trọng số chú ý trực tiếp từ các đặc
trưng điểm ảnh phong phú của\(\mathcal{\ P}_{w2}\) thay vì các đặc trưng
lọc thô của lớp tích chập chiều sâu:
\end{quote}

\[\mathcal{T}_{PDP} = c_{k} \otimes f_{k},\quad k = 1,\ldots,C_{out}\quad\left\lbrack 108,\text{Eq.~7} \right\rbrack\]

\begin{quote}
Trong đó, \(\otimes\) ký hiệu phép nhân theo từng kênh đặc trưng, và
\(f_{k} \in \mathcal{F}_{\mathcal{P}_{w2}}\) đại diện cho bộ trọng số
chú ý kênh được học từ \(\mathcal{P}_{w2}\).

\textbf{Cơ chế liên kết tàn dư đầy đủ.} Để duy trì liên kết đồng nhất
ngay cả khi kích thước không gian hoặc số lượng kênh của \(T\) thay đổi
so với đầu ra \(\ \mathcal{T}_{PDP}\) , TickNets đề xuất giải pháp sử
dụng một phép tích chập điểm phụ trợ \(\mathcal{P}_{wRs}\) để chuyển đổi
tensor đầu vào \(\mathcal{T}\) về cùng cấu trúc không gian và số lượng
kênh đặc trưng với đầu ra :
\end{quote}

\[\mathcal{P}_{wRs} = \nu_{j}*\mathcal{T},\quad j = 1,\ldots,C_{out}\quad\left\lbrack 110,\text{Eq.~8} \right\rbrack\]

\begin{quote}
Trong đó,
\(\mathcal{P}_{wRs} \in \mathbb{R}^{C_{out} \times H_{s} \times W_{s}}\)
sử dụng bước trượt \(s\) đồng bộ với bước trượt của lớp tích chập chiều
sâu \(\mathcal{D}_{ws}\) . Công thức tổng hợp đầu ra khối FR-PDP được
định nghĩa hoàn chỉnh như sau:
\end{quote}

\[\mathcal{T}_{FR - PDP} = \left\{ \begin{matrix}
\mathcal{T} \oplus \mathcal{T}_{PDP}, & \text{nếu~}s = 1\text{~và~}\left| \mathcal{T} \right| = \left| \mathcal{T}_{PDP} \right| \\
\mathcal{P}_{wRs} \oplus \mathcal{T}_{PDP}, & \text{trong~các~trường~hợp~khác} \\
\end{matrix} \right.\ \quad\left\lbrack 111,\text{Eq.~9} \right\rbrack\]

\begin{quote}
Trong đó, \(\oplus\) là phép toán cộng ma trận theo từng phần tử. Giải
pháp này đảm bảo mạng luôn khai thác được hai luồng thông tin tàn dư
song song để tối đa hóa hiệu năng trích xuất đặc trưng.
\end{quote}

\hypertarget{cux1ea5u-truxfac-xux1b0ux1a1ng-sux1ed1ng-co-giuxe3n-kuxeanh-huxecnh-dux1ea5u-tuxedch-vuxe0-muxf4-huxecnh-ticknet-small}{%
\subsubsection{Cấu trúc xương sống co giãn kênh hình dấu tích và mô hình
TickNet-small}\label{cux1ea5u-truxfac-xux1b0ux1a1ng-sux1ed1ng-co-giuxe3n-kuxeanh-huxecnh-dux1ea5u-tuxedch-vuxe0-muxf4-huxecnh-ticknet-small}}

\begin{quote}
Một đặc trưng độc bản của họ mạng TickNets là thiết kế hệ thống phân bổ
kênh đặc trưng theo mô hình co giãn kênh hình dấu tích phá vỡ tư duy
tăng kênh tuyến tính truyền thống để giảm thiểu tối đa sự bùng nổ tham
số thừa.

\textbf{Xương sống hình dấu tích cơ bản .} Một xương sống hình dấu tích
cơ bản được thiết lập bởi một chuỗi tuần tự gồm 5 khối FR-PDP. Sự co
giãn số lượng kênh đặc trưng đầu ra \(C_{out\ }\)của chuỗi 5 khối này
biến thiên tăng giảm nhịp nhàng mô phỏng theo hình dáng của một dấu tích
:
\end{quote}

\[Tick\_ B\left( \mathcal{T} \right) = \Psi_{s}^{512}\left( \Psi_{s}^{256}\left( \Psi_{s}^{128}\left( \Psi_{s}^{64}\left( \Psi_{s}^{128}\left( \mathcal{T} \right) \right) \right) \right) \right)\quad\left\lbrack 116,\text{Eq.~10} \right\rbrack\]

\begin{quote}
Trong đó, \(\Psi_{s}^{C_{out}}\ \)ký hiệu phép xử lý của một khối FR-PDP
cho ra số kênh \(C_{out}\) với bước trượt \(s\). Số lượng kênh đặc trưng
trải qua chu kỳ co giãn tuần tự:
\(128\  \rightarrow 64 \rightarrow 128 \rightarrow 256 \rightarrow 512\).

Việc co hẹp số lượng kênh đặc trưng về mức cực tiểu 64 kênh ở giai đoạn
đầu và chỉ mở rộng mạnh mẽ ở khối cuối cùng giúp kiểm soát nghiêm ngặt
dung lượng mô hình và loại bỏ các đặc trưng thừa.

\textbf{Kiến trúc mạng TickNet-small .} Tận dụng tính chất co giãn kênh
linh hoạt của xương sống cơ bản, mạng TickNet-small được xây dựng bằng
phương pháp ghép nối trực tiếp hai xương sống hình dấu tích tuần tự
thông qua cơ chế hooking:
\end{quote}

\[\mathcal{T}ick\mathcal{N}et\text{-small}\left( \mathcal{T} \right) = \mathcal{T}ick\_ B\left( \mathcal{T}ick\_ B\left( \mathcal{T} \right) \right)\quad\left\lbrack 119,\text{Eq.~11} \right\rbrack\]

\begin{quote}
Cấu trúc hoàn chỉnh của TickNet-small bao gồm 10 khối FR-PDP liên kết
chặt chẽ. Cấu hình phân bổ số lượng kênh đặc trưng cho 10 khối này được
triển khai chi tiết theo chuỗi co giãn đôi tương ứng với hai xương sống
hình dấu tích được ghép nối tiếp:
\(128 \rightarrow 64 \rightarrow 128 \rightarrow 256 \rightarrow 512 \rightarrow 128 \rightarrow 64 \rightarrow 128 \rightarrow 256 \rightarrow 512\).

Nhờ giải pháp co giãn kênh đặc trưng này, TickNet-small chỉ sở hữu vỏn
vẹn khoảng 3 triệu tham số huấn luyện trên tập dữ liệu ImageNet-1k. Con
số này tối giản hơn đáng kể so với mức 3.51 triệu tham số của mô hình
MobileNetV2 nhưng vẫn giúp mạng đạt được độ sâu trích xuất đặc trưng
vượt trội.
\end{quote}

\hypertarget{cux1a1-chux1ebf-lan-truyux1ec1n-ux111ux1eb7c-trux1b0ng-vuxe0-cuxe1ch-thux1ee9c-houx1ea1t-ux111ux1ed9ng-cux1ee7a-ticknets}{%
\subsubsection{Cơ chế lan truyền đặc trưng và cách thức hoạt động của
TickNets}\label{cux1a1-chux1ebf-lan-truyux1ec1n-ux111ux1eb7c-trux1b0ng-vuxe0-cuxe1ch-thux1ee9c-houx1ea1t-ux111ux1ed9ng-cux1ee7a-ticknets}}

\begin{quote}
Quy trình truyền xuôi (feed-forward pipeline) của mạng TickNet-small
hoạt động tuần tự để trích xuất đặc trưng hình ảnh phân cấp thông qua
các lớp tính toán được tối ưu hóa như sau:

{[}Ảnh đầu vào (224x224x3){]}

│

▼ (s=2)

{[}Lớp Initial\_conv3x3{]} ──► (112x112x32)

│

▼ (s=2)

{[}Backbone 1: Khối FR-PDP 1 - 5{]} ──► (7x7x512)

│

▼ (s=1)

{[}Điểm nối Hooking Point{]} ──► (CBAM cải tiến đề xuất)

│

▼ (s=1)

{[}Backbone 2: Khối FR-PDP 6 - 10{]} ──► (7x7x512)

│

▼ (s=1)

{[}Lớp Final\_conv1x1{]} ──► (7x7x1024)

│

▼

{[}Lớp Global AvgPooling{]} ──► (1024)

│

▼

{[}Lớp Fully Connected / Softmax{]} ──► {[}Nhãn dự đoán phân loại{]}

Giai đoạn xử lý ban đầu. Hình ảnh đầu vào kích thước chuẩn
\(224 \times 224 \times 3\) đi qua lớp tích chập tiêu chuẩn đầu tiên
\(Initial\_ conv3 \times 3\) có kích thước bộ lọc \(3 \times 3\) với
bước trượt \(s = 2\) để nén kích thước không gian xuống
\(112 \times 112\ \)và tăng kênh lên 32.

Giai đoạn lan truyền qua Backbone 1. Tensor đặc trưng được lan truyền
qua xương sống thứ nhất. Quá trình giảm kích thước bản đồ đặc trưng được
thực hiện một cách chiến lược bằng cách áp dụng bước trượt
\(s = 2\ \)tại khối 1, 3, 4, và 5. Đầu ra của backbone thứ nhất đạt kích
thước \(7 \times 7\) với số kênh đặc trưng là 512.

Giai đoạn lan truyền qua Backbone 2. Sau khi đi qua điểm nối mạng
(Hooking point), đặc trưng tiếp tục lan truyền qua xương sống dấu tích
thứ hai với bước trượt \(s = 1\) tại tất cả các khối để tập trung học
sâu các đặc trưng ngữ nghĩa nâng cao. Đầu ra cuối cùng duy trì kích
thước \(7 \times 7 \times 512\) .

Giai đoạn phân loại cuối. Bản đồ đặc trưng được tăng cường kênh đặc
trưng lên 1024 bằng lớp tích chập điểm \(Final\_ conv1 \times 1\), sau
đó làm phẳng qua lớp AvgPooling thành vector đặc trưng 1D có kích thước
1024, trước khi đi qua lớp kết nối đầy đủ và hàm kích hoạt Softmax để
tính toán xác suất phân loại lớp tương ứng.

Lưu ý cấu hình thực nghiệm: Đối với việc phân loại các hình ảnh kích
thước rất nhỏ như CIFAR-10 \(32 \times 32\), mạng TickNet-small sẽ cấu
hình bước trượt \(s = 1\ \)cho hai khối FR-PDP đầu tiên FR-PDP-1 và
FR-PDP-2 để bảo toàn tối đa các đặc trưng không gian mịn của ảnh tránh
bị mất mát thông tin quá sớm.
\end{quote}

\begin{enumerate}
\def\labelenumi{\arabic{enumi}.}
\setcounter{enumi}{1}
\item ~
  \hypertarget{cux1a1-chux1ebf-chuxfa-uxfd-attention-mechanism}{%
  \subsection{CƠ CHẾ CHÚ Ý (ATTENTION
  MECHANISM):}\label{cux1a1-chux1ebf-chuxfa-uxfd-attention-mechanism}}

  \begin{enumerate}
  \def\labelenumii{\arabic{enumii}.}
  \item ~
    \hypertarget{squeeze-and-excitation-mux1eb7c-ux111ux1ecbnh-trong-ticknets}{%
    \subsubsection{Squeeze-and-Excitation mặc định trong
    TickNets}\label{squeeze-and-excitation-mux1eb7c-ux111ux1ecbnh-trong-ticknets}}
  \end{enumerate}
\end{enumerate}

\begin{quote}
Cơ chế chú ý trong mạng nơ-ron nhân tạo là một kỹ thuật toán học mô
phỏng theo hệ thống nhận thức trực quan của con người, giúp mô hình tự
động gán trọng số và tập trung vào các vùng thông tin có giá trị cao
đồng thời bỏ qua các vùng thông tin nhiễu. Trong các kiến trúc mạng
nơ-ron tích chập siêu nhẹ, cơ chế chú ý theo kênh Squeeze-and-Excitation
(SE), được đề xuất bởi Hu và các cộng sự , là một giải pháp tối giản
nhưng mang lại hiệu năng biểu diễn đặc trưng vượt trội bằng cách thiết
lập các mối quan hệ phụ thuộc lẫn nhau giữa các kênh đặc trưng.

Khác biệt cốt lõi của việc tích hợp SE trong khối cơ bản FR-PDP của mạng
TickNets so với các mạng nơ-ron truyền thống nằm ở vị trí thu thập thông
tin. Thay vì tính toán trọng số chú ý từ các đặc trưng thô sau lớp tích
chập chiều sâu vốn chỉ mang tính chất lọc không gian đơn lẻ và chưa
thiết lập tương tác kênh TickNets thực hiện trích xuất thông tin chú ý
trực tiếp từ bản đồ đặc trưng điểm tích lũy \(\mathcal{P}_{w2}\) thu
được từ lớp tích chập điểm thứ hai :
\end{quote}

\[\mathcal{F}_{\mathcal{P}_{w2}} = \mathbf{F}_{ex}\left( \mathbf{F}_{sq}\left( \mathcal{P}_{w2} \right),\mathbf{W} \right)\]

\begin{quote}
Quy trình tính toán của mô-đun SE mặc định trong khối FR-PDP được thực
hiện qua hai giai đoạn toán học tuần tự:

\textbf{Phép toán Squeeze} (Nén đặc trưng toàn cục): Sử dụng phép gộp
trung bình toàn cục để nén thông tin không gian kích thước
\(H \times W\) của mỗi kênh đặc trưng thành một trị số vô hướng duy nhất
đại diện cho phân bố kênh toàn cục. Đầu ra của bước này là một vector
đặc trưng \(\mathbf{z} \in \mathbb{R}^{C_{out} \times 1 \times 1}\) với
phần tử thứ\(\ k\) được định nghĩa là :
\end{quote}

\[z_{k} = \frac{1}{H \times W}\sum_{i = 1}^{H}{\sum_{j = 1}^{W}\mathcal{P}_{w2}}(i,j,k)\]

\begin{quote}
\textbf{Phép toán Excitation} (Kích hoạt thích ứng): Để thu thập đầy đủ
các mối quan hệ phi tuyến giữa các kênh đặc trưng, vector nén \(z\)được
chuyển qua một cấu trúc Perceptron đa tầng tối giản với hai lớp kết nối
đầy đủ và một tỷ lệ giảm số kênh mặc định là \(r = 16\ \)để kiểm soát số
lượng tham số tính toán :
\end{quote}

\[\mathbf{s} = \sigma\left( \mathbf{W}_{2} \cdot \delta\left( \mathbf{W}_{1} \cdot \mathbf{z} \right) \right)\]

\begin{quote}
Trong đó, \(\delta\) ký hiệu hàm kích hoạt phi tuyến ReLU, \(\sigma\) là
hàm kích hoạt Sigmoid đưa giá trị trọng số về khoảng 1 và
\(\mathbf{W}_{1} \in \mathbb{R}^{\frac{C_{out}}{r} \times C_{out}}\) ,
\(\mathbf{W}_{2} \in \mathbb{R}^{C_{out} \times \frac{C_{out}}{r}}\) lần
lượt là các ma trận trọng số của hai lớp FC.

Cuối cùng, bản đồ đặc trưng đầu ra được hiệu chỉnh bằng cách nhân trực
tiếp từng phần tử của bản đồ đặc trưng điểm \(\mathcal{P}_{w2}\) với
trọng số chú ý kênh tương ứng \(s_{k},\ trong\ vecto\ \mathbf{s}\) :
\end{quote}

\[\mathcal{T}_{PDP} = c_{k} \otimes s_{k},\quad k = 1,\ldots,C_{out}\]

\hypertarget{channel-attention-module-cam}{%
\subsubsection{Channel Attention Module
(CAM):}\label{channel-attention-module-cam}}

\begin{quote}
\includegraphics[width=6.25197in,height=2.3937in]{vertopal_fe3807bd2608410aa17783721e5463e8/media/image3.PNG}Mô-đun
chú ý theo kênh Channel Attention Module trong cấu trúc chú ý hỗn hợp
CBAM đóng vai trò xác định thực thể thông tin nào mang tính quyết định
trong bản đồ đặc trưng đầu vào. Khác với cơ chế SE chỉ sử dụng duy nhất
phép gộp trung bình toàn cục GAP, CAM đề xuất một giải pháp thu thập
thông tin kênh song song bằng cách kết hợp đồng thời cả phép gộp trung
bình toàn cục GAP và phép gộp cực đại toàn cục Global Max Pooling.

Hình 2.2.2 Mô-đun chú ý kênh CAM

Giả sử tensor đặc trưng đầu vào tại một điểm nối mạng bất kỳ của TickNet
small là \(\mathbf{F} \in \mathbb{R}^{C \times H \times W}\) . Quy trình
tính toán của CAM được biểu diễn thông qua các bước toán học chặt chẽ
sau:

\textbf{Trích xuất đặc trưng kênh song song}: Bản đồ đặc trưng
\(\mathbf{F}\) được nén độc lập theo chiều không gian để tạo ra hai
vector đặc trưng kênh riêng biệt đại diện cho hai khía cạnh thông tin
khác nhau:

Vector trung bình toàn cục:
\(\mathbf{F}_{avg}^{c} \in \mathbb{R}^{C \times 1 \times 1}\ \)lưu trữ
thông tin bối cảnh không gian mang tính trung bình.

Vecto cực đại toàn cục:
\(\mathbf{F}_{\max}^{c} \in \mathbb{R}^{C \times 1 \times 1}\) định vị
các đặc trưng nổi bật nhất của đối tượng phân loại trong không gian.

\textbf{Lan truyền qua mạng Perceptron đa tầng chia sẻ} (Shared MLP): Để
đảm bảo tính tối giản về mặt tham số huấn luyện của mạng siêu nhẹ, cả
hai vector đặc trưng trên được đưa qua một cấu trúc MLP chia sẻ trọng số
duy nhất. Shared MLP này bao gồm một lớp ẩn với số lượng nơ-ron được
giảm đi theo tỷ lệ giảm \(r\) với hệ số nén mặc định \(r = 16\) hoặc
điều chỉnh linh hoạt tùy thuộc cấu hình thực nghiệm:
\end{quote}

\[\mathbf{M}_{c}\left( \mathbf{F} \right) = \sigma\left( \text{MLP}\left( \text{AvgPool}\left( \mathbf{F} \right) \right) + \text{MLP}\left( \text{MaxPool}\left( \mathbf{F} \right) \right) \right)\]

\[\mathbf{M}_{\mathbf{c}}\left( \mathbf{F} \right)\mathbf{= \sigma}\left( \mathbf{W}_{\mathbf{1}}\mathbf{\cdot}\left( \mathbf{W}_{\mathbf{0}}\mathbf{\cdot}\mathbf{F}_{\mathbf{avg}}^{\mathbf{c}} \right)\mathbf{+}\mathbf{W}_{\mathbf{1}}\mathbf{\cdot}\left( \mathbf{W}_{\mathbf{0}}\mathbf{\cdot}\mathbf{F}_{\mathbf{\max}}^{\mathbf{c}} \right) \right)\]

\begin{quote}
Trong đó, \(\mathbf{W}_{0} \in \mathbb{R}^{\frac{C}{r} \times C}\) và
\(\mathbf{W}_{1} \in \mathbb{R}^{C \times \frac{C}{r}}\) là các trọng số
liên kết dùng chung cho cả hai nhánh xử lý thông tin, và \(\sigma\)đại
diện cho hàm kích hoạt Sigmoid. Bản đồ chú ý kênh cuối cùng thu được là
một vector trọng số phân bố thích ứng
\(\mathbf{M}_{c}\left( \mathbf{F} \right) \in \mathbb{R}^{C \times 1 \times 1}\).
\end{quote}

\hypertarget{spatial-attention-module-sam}{%
\subsubsection{Spatial Attention Module
(SAM):}\label{spatial-attention-module-sam}}

\begin{quote}
\includegraphics[width=6.25in,height=2.08333in]{vertopal_fe3807bd2608410aa17783721e5463e8/media/image4.png}Trong
khi mô-đun CAM tập trung vào việc tìm kiếm các đặc trưng kênh ý nghĩa,
mô-đun chú ý theo không gian Spatial Attention Module (SAM) chịu trách
nhiệm định vị vị trí không gian chứa thông tin cốt lõi cần được chú ý
trong ảnh. SAM đóng vai trò là một bộ lọc không gian cục bộ mạnh mẽ,
giúp loại bỏ nhiễu từ bối cảnh phức tạp của môi trường xung quanh đối
tượng phân loại.

Hình 2.2.3 Sơ đồ mạng nơ-ron Tensor

Quy trình thiết lập bản đồ chú ý không gian 2D từ một tensor đặc trưng
đầu vào \(\mathbf{F}' \in \mathbb{R}^{C \times H \times W}\) được thực
hiện như sau:

\textbf{Gộp đặc trưng dọc theo chiều sâu kênh}: SAM thực hiện phép toán
gộp trung bình (Average Pooling) và gộp cực đại (Max Pooling) dọc theo
trục kênh để nén thông tin chiều sâu và trích xuất trực tiếp mối quan hệ
không gian 2D:

Bản đồ trung bình không gian:
\(\mathbf{F}_{avg}^{s} \in \mathbb{R}^{1 \times H \times W}\)

Bản đồ cực đại không gian:
\(\mathbf{F}_{\max}^{s} \in \mathbb{R}^{1 \times H \times W}\)

Ghép nối và tích chập tổng hợp thông tin: Hai bản đồ đặc trưng 2D trên
được ghép nối dọc theo chiều sâu để tạo thành một tensor đặc trưng hỗn
hợp kích thước \(\mathbb{R}^{2 \times H \times W}\) . Tiếp theo, một
phép toán tích chập tiêu chuẩn với kích thước bộ lọc lớn \(k \times k\)
thường mặc định \(7\  \times \ 7\ \)đối với ảnh lớn hoặc tối ưu hóa
xuống \(3 \times 3\ \)đối với ảnh nhỏ được áp dụng để trích xuất cấu
trúc không gian diện rộng:
\end{quote}

\[\mathbf{M}_{s}\left( \mathbf{F}' \right) = \sigma\left( f^{k \times k}\left( \left\lbrack \text{AvgPool}\left( \mathbf{F}' \right);\text{MaxPool}\left( \mathbf{F}' \right) \right\rbrack \right) \right)\]

\[\mathbf{M}_{s}\left( \mathbf{F}' \right) = \sigma\left( f^{k \times k}\left( \left\lbrack \mathbf{F}_{avg}^{s};\mathbf{F}_{\max}^{s} \right\rbrack \right) \right)\]

\begin{quote}
Trong đó, \(f^{k \times k}\) đại diện cho phép toán tích chập với bộ lọc
kích thước \(k \times k\)và \(\sigma\) diện cho hàm kích hoạt Sigmoid
đưa phân bố không gian về khoảng giá trị tối ưu. Kết quả thu được là bản
đồ chú ý không gian 2D
\(\mathbf{M}_{s}\left( \mathbf{F}' \right) \in \mathbb{R}^{1 \times H \times W}\).
\end{quote}

\hypertarget{muxf4-ux111un-chuxfa-uxfd-khux1ed1i-tuxedch-chux1eadp-cbam-houx1ea1t-ux111ux1ed9ng-tuux1ea7n-tux1ef1-cam-sam}{%
\subsubsection{Mô-đun chú ý khối tích chập CBAM hoạt động tuần tự CAM +
SAM}\label{muxf4-ux111un-chuxfa-uxfd-khux1ed1i-tuxedch-chux1eadp-cbam-houx1ea1t-ux111ux1ed9ng-tuux1ea7n-tux1ef1-cam-sam}}

\begin{quote}
CBAM là sự kết hợp đồng bộ và tuần tự của hai mô-đun CAM và SAM nhằm tối
đa hóa khả năng biểu diễn đặc trưng phân cấp của mạng nơ-ron tích chập.
Khác với các phương pháp chèn song song hoặc chèn độc lập, CBAM thiết
lập một quy trình tinh lọc thông tin hai bước cực kỳ khoa học: thông tin
đi qua bộ lọc kênh CAM trước để xác định các kênh đặc trưng quan trọng,
sau đó đi qua bộ lọc không gian SAM để khoanh vùng tọa độ đặc trưng mục
tiêu.

Sự phối hợp tuần tự này được định nghĩa một cách chặt chẽ thông qua hệ
thống công thức sau:
\end{quote}

\[\mathbf{F}' = \mathbf{M}_{c}\left( \mathbf{F} \right) \otimes \mathbf{F}\]

\[\mathbf{F}'' = \mathbf{M}_{s}\left( \mathbf{F}' \right) \otimes \mathbf{F}'\]

\begin{quote}
\includegraphics[width=5.32569in,height=2.71042in]{vertopal_fe3807bd2608410aa17783721e5463e8/media/image5.png}Trong
đó, \(F\) là tensor đặc trưng thô ban đầu, \(F'\) là bản đồ đặc trưng đã
được hiệu chỉnh thông tin theo kênh thông qua phép nhân element-wise
\(\otimes \ \)với vector chú ý
\(\mathbf{M}_{c}\left( \mathbf{F} \right)\) , và \(F''\) là bản đồ đặc
trưng đầu ra hoàn chỉnh sau khi được lọc không gian diện rộng bởi bản đồ
nhiệt chú ý 2D \(\mathbf{M}_{s}(\mathbf{F}')\).

Hình 2.2.4 Sơ đồ quy trình CBAM

Việc áp dụng cơ chế CBAM tuần tự mang lại ba lợi ích to lớn cho bài toán
phân loại hình ảnh:

\textbf{Tính toàn diện của biểu diễn đặc trưng:} Giúp mô hình đồng thời
nắm bắt cả thông tin mức độ cao hình dạng, nhãn lớp thông qua CAM và
thông tin không gian mức độ thấp góc, cạnh, định vị vật thể thông qua
SAM.

\textbf{Khả năng tương thích kiến trúc cực cao:} CBAM là một mô-đun độc
lập và cực kỳ gọn nhẹ chỉ bổ sung một lượng tham số không đáng kể từ
Shared MLP, cho phép tích hợp linh hoạt vào bất kỳ điểm nối mạng nào mà
không làm ảnh hưởng đến tiến trình truyền xuôi gốc.

\textbf{Tối ưu hóa khả năng giải thích của mô hình:} Các bản đồ chú ý
không gian 2D tạo ra bởi SAM cung cấp nền tảng toán học rõ ràng giúp
người nghiên cứu dễ dàng giải thích hành vi học máy trực quan bằng các
thuật toán bản đồ nhiệt như Grad-CAM.
\end{quote}

\hypertarget{phux1b0ux1a1ng-phuxe1p-nghiuxean-cux1ee9u}{%
\subsection{PHƯƠNG PHÁP NGHIÊN
CỨU}\label{phux1b0ux1a1ng-phuxe1p-nghiuxean-cux1ee9u}}

\begin{quote}
Đề tài được triển khai dựa trên sự kết hợp giữa phương pháp nghiên cứu
lý thuyết và phương pháp nghiên cứu thực nghiệm.

\textbf{Về phương pháp lý thuyết:} Đề tài tiến hành thu thập, phân tích
và hệ thống hóa các tài liệu khoa học liên quan đến mạng nơ-ron tích
chập siêu nhẹ, trọng tâm là họ mạng TickNets với khối cơ bản FR-PDP và
cấu trúc xương sống co giãn kênh hình dấu tích. Đồng thời, nghiên cứu cơ
chế chú ý theo kênh SE mặc định kết hợp với mô đun chú ý hỗn hợp tuần tự
CBAM (gồm CAM và SAM). Trên cơ sở đó, đề tài phân tích toán học để chỉ
ra sự chồng chéo không gian của attention cục bộ bên trong khối và đề
xuất giải pháp chèn CBAM phân cấp tại các điểm nối mạng để tối ưu hiệu
năng.

\textbf{Về phương pháp thực nghiệm:} Đề tài lập trình hiện thực hóa giải
pháp cải tiến trên ngôn ngữ Python với thư viện TensorFlow và Keras. Quy
trình thực nghiệm được triển khai đồng bộ qua 4 bước tuần tự sau:
\end{quote}

\begin{enumerate}
\def\labelenumi{\arabic{enumi}.}
\item
  \begin{quote}
  Chuẩn bị và tiền xử lý dữ liệu: Thử nghiệm trên 4 bộ dữ liệu, chia làm
  hai nhóm: nhóm học thưuật Fashion-MNIST, CIFAR-10 và nhóm thực tiễn
  (PlantVillage, Chest-Xray-Pneumonia). Pixel ảnh được chuẩn hóa về
  khoảng , nhãn được mã hóa One-Hot và áp dụng các kỹ thuật tăng cường
  dữ liệu (xoay, lật, dịch chuyển) để chống quá khớp.
  \end{quote}
\item
  Xây dựng các mô hình đối chứng: Thiết lập 3 kiến trúc mạng gồm: Mô
  hình 1 SE Baseline TickNet-small gốc sử dụng SE trong khối. Mô hình 2
  CBAM-Local chèn CBAM thay thế SE cục bộ. Mô hình 3 CBAM Hook đề xuất
  cải tiến chèn CBAM phân cấp tại điểm nối và cuối mạng.
\item
  Huấn luyện và tối ưu hóa: Huấn luyện các mô hình bằng bộ tối ưu hóa
  SGD kết hợp hàm loss tương ứng. Quá trình huấn luyện tích hợp bộ đôi
  hàm gọi lại động là tự động giảm tốc độ học và dừng huấn luyện sớm để
  khôi phục bộ trọng số tối ưu nhất.
\item
  Đánh giá và trực quan hóa giải thích: Đo lường hiệu năng các mô hình
  thông qua các độ đo chuẩn mực gồm Accuracy, Precision, Recall,
  F1-score và ma trận Confusion Matrix. Cuối cùng, ứng dụng thuật toán
  Grad-CAM để trực quan hóa bản đồ nhiệt, giải thích khoa học khả năng
  tập trung đặc trưng không gian của giải pháp cải tiến đề xuất.
\end{enumerate}

\hypertarget{chux1b0ux1a1ng-3-buxe0i-touxe1n-phuxe2n-loux1ea1i-huxecnh-ux1ea3nh-vuxe0-ux111ux1ec1-xuux1ea5t-cux1ea3i-tiux1ebfn}{%
\section{\texorpdfstring{CHƯƠNG 3\\
BÀI TOÁN PHÂN LOẠI HÌNH ẢNH VÀ ĐỀ XUẤT CẢI
TIẾN}{CHƯƠNG 3 BÀI TOÁN PHÂN LOẠI HÌNH ẢNH VÀ ĐỀ XUẤT CẢI TIẾN}}\label{chux1b0ux1a1ng-3-buxe0i-touxe1n-phuxe2n-loux1ea1i-huxecnh-ux1ea3nh-vuxe0-ux111ux1ec1-xuux1ea5t-cux1ea3i-tiux1ebfn}}

\begin{enumerate}
\def\labelenumi{\arabic{enumi}.}
\item ~
  \hypertarget{giux1edbi-thiux1ec7u-buxe0i-touxe1n-phuxe2n-loux1ea1i-huxecnh-ux1ea3nh}{%
  \subsection{GIỚI THIỆU BÀI TOÁN PHÂN LOẠI HÌNH
  ẢNH}\label{giux1edbi-thiux1ec7u-buxe0i-touxe1n-phuxe2n-loux1ea1i-huxecnh-ux1ea3nh}}

  \begin{enumerate}
  \def\labelenumii{\arabic{enumii}.}
  \item ~
    \hypertarget{phuxe1t-biux1ec3u-buxe0i-touxe1n}{%
    \subsubsection{Phát biểu bài
    toán}\label{phuxe1t-biux1ec3u-buxe0i-touxe1n}}
  \end{enumerate}
\end{enumerate}

\begin{quote}
Phân loại hình ảnh là một trong những bài toán nền tảng và quan trọng
nhất trong lĩnh vực thị giác máy tính. Bài toán được phát biểu một cách
hình thức như sau: cho một hình ảnh đầu vào \emph{I} và một tập hữu hạn
các nhãn lớp
\(\mathcal{C} = \left\{ c_{1},c_{2},\ldots,c_{K} \right\}\ \)đã được xác
định trước, hệ thống cần xây dựng một hàm ánh xạ
\(f:I \rightarrow \mathcal{C\ }\)để gán nhãn lớp phù hợp nhất cho hình
ảnh đó. Nói cách khác, mục tiêu là xác định hình ảnh thuộc vào nhóm nào
trong số các nhóm đã được định nghĩa sẵn.

Ví dụ minh hoạ cụ thể: trong một hệ thống nhận dạng động vật, tập nhãn
có thể bao gồm
\(\mathcal{C} = \left\{ \text{chó},\text{mèo},\text{chim},\text{ngựa} \right\}.\ \)Khi
nhận một bức ảnh chụp một con chó Golden Retriever, hệ thống cần phân
tích các đặc trưng thị giác của ảnh bao gồm hình dạng, kết cấu lông, tỉ
lệ cơ thể để đưa ra dự đoán chính xác rằng ảnh thuộc lớp ``chó'' với xác
suất tin cậy cao nhất.

\includegraphics[width=4.99925in,height=3.29264in]{vertopal_fe3807bd2608410aa17783721e5463e8/media/image6.png}

\emph{Hình 3.1.1a Tổng quan bài toán phân loại ảnh}
\end{quote}

\hypertarget{biux1ec3u-diux1ec5n-huxecnh-ux1ea3nh-trong-muxe1y-tuxednh}{%
\subsubsection{\texorpdfstring{Biểu diễn hình ảnh trong máy tính
}{Biểu diễn hình ảnh trong máy tính }}\label{biux1ec3u-diux1ec5n-huxecnh-ux1ea3nh-trong-muxe1y-tuxednh}}

\begin{quote}
Đối với con người, việc nhận dạng nội dung một bức ảnh diễn ra gần như
tức thì và tự nhiên. Tuy nhiên, đối với máy tính, một hình ảnh chỉ đơn
thuần là một mảng ba chiều chứa các giá trị số nguyên. Cụ thể, một ảnh
màu RGB được biểu diễn dưới dạng tensor
\(\mathbf{X} \in \mathbb{R}^{H \times W \times C}\), trong đó \(H\) là
chiều cao, \(W\ \)là chiều rộng tính bằng số pixel, và \(C = 3\ \)tương
ứng với ba kênh màu đỏ , xanh lá , xanh dương. Mỗi phần tử trong tensor
nhận giá trị từ 0 đến 255, thể hiện cường độ sáng của pixel tại vị trí
tương ứng trên kênh màu đó.

Ví dụ, một bức ảnh kích thước \(224 \times 224\) pixel với 3 kênh màu sẽ
chứa tổng cộng \(224 \times 224 \times 3 = 150.528\ \)giá trị số. Nhiệm
vụ của thuật toán phân loại chính là từ tập hợp hàng trăm nghìn con số
này, rút trích ra được các đặc trưng có ý nghĩa ngữ nghĩa để phân biệt
giữa các lớp đối tượng khác nhau. Khoảng cách lớn giữa thông tin mức
pixel thô và ý nghĩa ngữ nghĩa cấp cao này được gọi là khoảng trống ngữ
nghĩa, và đây chính là thách thức cốt lõi mà các phương pháp phân loại
ảnh cần giải quyết.

\includegraphics[width=3.41293in,height=3.53788in]{vertopal_fe3807bd2608410aa17783721e5463e8/media/image7.jpeg}

\emph{Hình 3.1.2 So sánh cách con người nhìn ảnh trực quan và cách máy
tính biểu diễn ảnh tensor} \(H \times W \times C\ \)\emph{các giá trị
pixel}
\end{quote}

\hypertarget{thuxe1ch-thux1ee9c-cux1ee7a-buxe0i-touxe1n}{%
\subsubsection{Thách thức của bài
toán}\label{thuxe1ch-thux1ee9c-cux1ee7a-buxe0i-touxe1n}}

Mặc dù bài toán phân loại hình ảnh có thể được phát biểu một cách đơn
giản, việc xây dựng một hệ thống phân loại đạt độ chính xác cao trong
thực tế đối mặt với nhiều thách thức nghiêm trọng:

\includegraphics[width=3.41333in,height=3.41333in]{vertopal_fe3807bd2608410aa17783721e5463e8/media/image8.jpeg}

\emph{Hình 3.1.3 Bốn thách thức chính của bài toán phân loại ảnh: biến
đổi góc nhìn, biến đổi tỉ lệ, che khuất và đa dạng nội lớp}

\textbf{Biến đổi góc nhìn (Viewpoint Variation):} Cùng một đối tượng khi
được chụp từ các góc độ khác nhau sẽ tạo ra các hình ảnh có hình dạng
rất khác biệt trên mặt phẳng 2D. Ví dụ, một chiếc ô tô chụp từ phía
trước trông hoàn toàn khác so với khi chụp từ phía trên.

\textbf{Biến đổi tỉ lệ (Scale Variation):} Đối tượng có thể xuất hiện ở
nhiều kích thước khác nhau trong ảnh, từ rất nhỏ chiếm vài pixel đến rất
lớn phủ gần toàn bộ khung hình. Hệ thống phải nhận dạng được đối tượng
bất kể tỉ lệ.

\textbf{Hiện tượng che khuất (Occlusion):} Trong nhiều tình huống thực
tế, đối tượng cần nhận dạng bị che khuất một phần hoặc phần lớn bởi các
vật thể khác. Khi chỉ một phần nhỏ của đối tượng hiển thị, việc phân
loại trở nên rất khó khăn.

\textbf{Đa dạng nội lớp (Intra-class Variation):} Các đối tượng thuộc
cùng một lớp có thể có hình dạng, màu sắc và kết cấu rất khác nhau.
Chẳng hạn, lớp ``ghế'' bao gồm ghế sofa, ghế gỗ, ghế xoay văn phòng tất
cả đều là ghế nhưng có hình dáng bên ngoài rất khác biệt.

\textbf{Biến đổi điều kiện chiếu sáng (Illumination Variation):} Cùng
một đối tượng dưới các điều kiện ánh sáng khác nhau ngoài trời, trong
nhà, ban đêm sẽ tạo ra các giá trị pixel rất khác nhau, gây khó khăn cho
việc nhận dạng.

\textbf{Nhiễu nền phức tạp (Background Clutter):} Khi đối tượng có màu
sắc hoặc kết cấu tương tự với nền xung quanh, ranh giới giữa đối tượng
và nền trở nên mờ nhạt, khiến việc tách biệt và phân loại gặp nhiều trở
ngại.

\hypertarget{ux1ee9ng-dux1ee5ng-thux1ef1c-tiux1ec5n}{%
\subsubsection{Ứng dụng thực
tiễn}\label{ux1ee9ng-dux1ee5ng-thux1ef1c-tiux1ec5n}}

\begin{quote}
Phân loại hình ảnh ngày nay đóng vai trò then chốt trong rất nhiều lĩnh
vực ứng dụng thực tiễn, mang lại giá trị to lớn cho cả đời sống và sản
xuất:

\textbf{Y tế và chẩn đoán hình ảnh:} Phân loại ảnh X-quang để phát hiện
viêm phổi, ung thư phổi, phân tích ảnh mô bệnh học để hỗ trợ bác sĩ chẩn
đoán chính xác hơn, sàng lọc bệnh lý võng mạc từ ảnh quang học đáy mắt.
Đặc biệt, bộ dữ liệu Chest-Xray-Pneumonia được sử dụng trong nghiên cứu
này chính là một ví dụ điển hình cho ứng dụng phân loại ảnh y tế.

\textbf{Nông nghiệp thông minh:} Nhận dạng và phân loại bệnh trên lá cây
trồng từ ảnh chụp, giúp nông dân phát hiện sớm dịch bệnh và có biện pháp
xử lý kịp thời. Bộ dữ liệu PlantVillage trong tiểu luận này là minh
chứng cho tiềm năng ứng dụng của công nghệ phân loại ảnh trong nông
nghiệp chính xác.

\textbf{Xe tự hành:} Nhận dạng biển báo giao thông, phân loại đối tượng
trên đường người đi bộ, xe cộ, chướng ngại vật là thành phần không thể
thiếu trong hệ thống lái xe tự động.

\textbf{Kiểm tra chất lượng sản phẩm:} Trong các dây chuyền sản xuất
công nghiệp, hệ thống phân loại ảnh tự động phát hiện sản phẩm lỗi, kiểm
tra khuyết tật bề mặt với tốc độ và độ chính xác vượt trội so với kiểm
tra thủ công.

\textbf{An ninh và giám sát:} Nhận dạng khuôn mặt, phân loại hành vi bất
thường trong video giám sát, hỗ trợ lực lượng an ninh trong công tác
phòng chống tội phạm.

\textbf{Viễn thám và môi trường:} Phân loại ảnh vệ tinh để giám sát sử
dụng đất, theo dõi biến đổi rừng, dự báo thiên tai và đánh giá tác động
môi trường.
\end{quote}

\hypertarget{cuxe1c-phux1b0ux1a1ng-phuxe1p-tiux1ebfp-cux1eadn-buxe0i-touxe1n-phuxe2n-loux1ea1i-huxecnh-ux1ea3nh}{%
\subsubsection{Các phương pháp tiếp cận bài toán phân loại hình
ảnh}\label{cuxe1c-phux1b0ux1a1ng-phuxe1p-tiux1ebfp-cux1eadn-buxe0i-touxe1n-phuxe2n-loux1ea1i-huxecnh-ux1ea3nh}}

\begin{quote}
Lịch sử phát triển của bài toán phân loại ảnh có thể chia thành hai giai
đoạn chính:

\textbf{Giai đoạn 1 - Phương pháp truyền thống (trước 2012):} Các phương
pháp sử dụng kỹ thuật trích xuất đặc trưng thủ công như SIFT,HOG hay
LBP, kết hợp với các bộ phân lớp cổ điển như SVM hoặc KNN. Hạn chế chính
của nhóm phương pháp này là đặc trưng được thiết kế bởi con người, phụ
thuộc nhiều vào chuyên gia và không có khả năng tự thích ứng với dữ liệu
mới.

\textbf{Giai đoạn 2 - Phương pháp học sâu (từ 2012 đến nay):} Bước ngoặt
lớn đến từ năm 2012 khi mô hình AlexNet của Krizhevsky và cộng sự giành
chiến thắng cuộc thi ImageNet Large Scale Visual Recognition Challenge
với khoảng cách vượt trội so với các phương pháp truyền thống. Kể từ đó,
mạng nơ-ron tích chập trở thành kiến trúc chủ đạo cho bài toán phân loại
ảnh. Các kiến trúc CNN đặc biệt thành công bao gồm VGGNet (2014),
GoogLeNet/Inception (2014), ResNet (2015), DenseNet (2017), và
EfficientNet (2019). Ưu điểm cốt lõi của phương pháp học sâu là khả năng
tự động học các đặc trưng phân cấp trực tiếp từ dữ liệu thô, mà không
cần sự can thiệp thủ công trong quá trình thiết kế đặc trưng.

Tuy nhiên, các mô hình CNN hiệu năng cao thường có kích thước rất lớn
hàng chục đến hàng trăm triệu tham số, đòi hỏi tài nguyên tính toán mạnh
mẽ, không phù hợp với các thiết bị có tài nguyên hạn chế như điện thoại
di động, thiết bị IoT hay thiết bị bay không người lái. Chính nhu cầu
cân bằng giữa hiệu năng phân loại và chi phí tính toán đã thúc đẩy sự
phát triển của các kiến trúc mạng CNN siêu nhẹ light-weight CNNs như
MobileNets, ShuffleNets và TickNets là đối tượng nghiên cứu chính của
tiểu luận này.
\end{quote}

\hypertarget{vai-truxf2-cux1ee7a-cux1a1-chux1ebf-chuxfa-uxfd-trong-buxe0i-touxe1n-phuxe2n-loux1ea1i-ux1ea3nh}{%
\subsubsection{Vai trò của cơ chế chú ý trong bài toán phân loại
ảnh}\label{vai-truxf2-cux1ee7a-cux1a1-chux1ebf-chuxfa-uxfd-trong-buxe0i-touxe1n-phuxe2n-loux1ea1i-ux1ea3nh}}

\begin{quote}
Bên cạnh việc tối ưu kiến trúc mạng, một hướng nghiên cứu quan trọng
khác nhằm nâng cao hiệu năng phân loại ảnh là tích hợp cơ chế chú ý . Cơ
chế chú ý mô phỏng khả năng tập trung chọn lọc của thị giác con người:
thay vì xử lý đồng đều toàn bộ hình ảnh, mô hình học cách nhấn mạnh các
vùng và đặc trưng quan trọng trong khi giảm thiểu ảnh hưởng của nhiễu
nền và thông tin không liên quan.

Đặc biệt trong các bài toán phân loại ảnh thực tế nơi đối tượng cần nhận
dạng thường bị lẫn trong nền phức tạp .Ví dụ: vùng bệnh nhỏ trên lá cây
giữa nền lá xanh, hay vùng viêm mờ trên ảnh X-quang phổi. Cơ chế chú ý
đóng vai trò quyết định trong việc giúp mô hình định vị chính xác vùng
đặc trưng mang tính phân biệt cao, từ đó cải thiện đáng kể độ chính xác
phân loại.

Chính từ bối cảnh trên, tiểu luận này đặt ra câu hỏi nghiên cứu: Làm thế
nào để tích hợp hiệu quả cơ chế chú ý CBAM vào kiến trúc mạng siêu nhẹ
TickNets nhằm nâng cao hiệu năng phân loại ảnh mà không làm bùng nổ chi
phí tính toán? Câu trả lời cho câu hỏi này sẽ được trình bày chi tiết
trong phần đề xuất mô hình cải tiến phân cấp ở mục 3.2 tiếp theo.
\end{quote}

\hypertarget{ux111ux1ec1-xuux1ea5t-muxf4-huxecnh-ticknets-kux1ebft-hux1ee3p-cux1a1-chux1ebf-chuxfa-uxfd-phuxe2n-cux1ea5p-hierarchical-attention}{%
\subsection{ĐỀ XUẤT MÔ HÌNH TICKNETS KẾT HỢP CƠ CHẾ CHÚ Ý PHÂN CẤP
HIERARCHICAL
ATTENTION}\label{ux111ux1ec1-xuux1ea5t-muxf4-huxecnh-ticknets-kux1ebft-hux1ee3p-cux1a1-chux1ebf-chuxfa-uxfd-phuxe2n-cux1ea5p-hierarchical-attention}}

\begin{quote}
Dựa trên nền tảng lý thuyết ở Chương 2 và bối cảnh bài toán ở mục 3.1,
mục này đề xuất kiến trúc cải tiến TickNet-small CBAM-Hook một hệ thống
chú ý phân cấp hai tầng giữ nguyên SE cục bộ bên trong mỗi khối FR-PDP,
đồng thời chèn CBAM toàn cục tại các điểm nối chiến lược hooking points
giữa hai nhánh xương sống và trước lớp GAP cuối mạng.
\end{quote}

\hypertarget{thiux1ebft-kux1ebf-kiux1ebfn-truxfac-ux111ux1ec1-xuux1ea5t-ticknet-small-cbam-hook}{%
\subsubsection{3.2.1 Thiết kế kiến trúc đề xuất TickNet-small
CBAM-Hook}\label{thiux1ebft-kux1ebf-kiux1ebfn-truxfac-ux111ux1ec1-xuux1ea5t-ticknet-small-cbam-hook}}

\begin{quote}
TickNet-small nguyên bản tích hợp SE trong mỗi khối FR-PDP để tái hiệu
chỉnh kênh cục bộ, nhưng hoàn toàn thiếu cơ chế lọc theo không gian
không xác định được vùng nào trên bản đồ đặc trưng chứa thông tin phân
loại quan trọng nhất. Hạn chế này đặc biệt nghiêm trọng trên dữ liệu
thực tế có nền phức tạp nơi vùng bệnh lý chỉ chiếm phần nhỏ diện tích
ảnh.

Nghiên cứu so sánh hai chiến lược chèn CBAM:

\textbf{CBAM-Local:} Thay trực tiếp SE bằng CBAM bên trong mỗi khối
FR-PDP. Cách tiếp cận trực quan nhưng tiềm ẩn xung đột kiến trúc với
phép tích chập tách biệt chiều sâu.

\textbf{CBAM-Hook (đề xuất):} Giữ nguyên SE nội khối, chèn CBAM ở bên
ngoài các khối tại các vị trí chiến lược, tạo hệ thống chú ý phân cấp
hai tầng bổ sung cho nhau.

Ba nguyên lý thiết kế cốt lõi:
\end{quote}

\begin{enumerate}
\def\labelenumi{\arabic{enumi}.}
\item
  Phân tách chức năng: SE nội khối lọc kênh nào quan trọng ở phạm vi cục
  bộ. CBAM ngoại khối lọc vùng không gian nào quan trọng trên bản đồ đặc
  trưng đã tích lũy qua nhiều khối.
\item
  Tránh chồng chéo đặc trưng: SAM của CBAM dùng gộp kênh để tạo bản đồ
  chú ý không gian 2D. Khi đặt SAM cùng mức với DWConv trong FR-PDP, cả
  hai cùng hoạt động trên cùng miền không gian cục bộ với receptive
  field tương đồng, gây dư thừa thông tin.
\item
  Chèn tại điểm chuyển pha: Tick-shape Backbone gồm pha mở rộng và pha
  thu hẹp kênh. Điểm chuyển tiếp nơi đặc trưng đạt mức trừu tượng cao
  nhất là vị trí lý tưởng để CBAM tinh lọc toàn cục trước khi đặc trưng
  bước vào pha nén.
\end{enumerate}

\begin{quote}
Ba điểm chèn chiến lược. Kiến trúc CBAM-Hook chèn CBAM tại ba hooking
points:

Hooking Point 1 - Giữa Expanding và Contracting Backbone\textbf{:} Vị
trí số kênh cực đại, CBAM tái hiệu chỉnh kênh và lọc không gian trước
khi đặc trưng vào pha nén.

Hooking Point 2 - Cuối pha thu hẹp: ``Cổng kiểm soát cuối
cùng\textquotesingle'' tinh chỉnh kênh và không gian của bản đồ đặc
trưng đã nén trước khi truyền tới lớp phân loại.

Hooking Point 3 - Trước GAP\textbf{:} Lọc không gian lần cuối, đảm bảo
GAP chỉ tính trung bình trên vùng chứa đối tượng, không bị pha loãng bởi
nền.

\includegraphics[width=4.2197in,height=4.2197in]{vertopal_fe3807bd2608410aa17783721e5463e8/media/image9.jpeg}\emph{Hình
3.2.1: Kiến trúc tổng thể mô hình đề xuất TickNet-small CBAM-Hook}
\end{quote}

\begin{longtable}[]{@{}
  >{\raggedright\arraybackslash}p{(\columnwidth - 6\tabcolsep) * \real{0.3774}}
  >{\raggedright\arraybackslash}p{(\columnwidth - 6\tabcolsep) * \real{0.1913}}
  >{\raggedright\arraybackslash}p{(\columnwidth - 6\tabcolsep) * \real{0.2105}}
  >{\raggedright\arraybackslash}p{(\columnwidth - 6\tabcolsep) * \real{0.2209}}@{}}
\toprule()
\begin{minipage}[b]{\linewidth}\raggedright
Thành phần kiến trúc
\end{minipage} & \begin{minipage}[b]{\linewidth}\raggedright
SE Baseline
\end{minipage} & \begin{minipage}[b]{\linewidth}\raggedright
CBAM-Local
\end{minipage} & \begin{minipage}[b]{\linewidth}\raggedright
CBAM-Hook
\end{minipage} \\
\midrule()
\endhead
Attention nội khối FR-PDP & SE & CBAM & SE giữ nguyên \\
Attention tại Hooking Point 1 & Không & Không & CBAM \\
Attention tại Hooking Point 2 & Không & Không & CBAM \\
Attention trước GAP & Không & Không & CBAM \\
Lọc kênh cục bộ & Có SE & Có CAM & Có SE \\
Lọc không gian cục bộ & Không & Có SAM & Không \\
Lọc không gian toàn cục & Không & Không & Có SAM \\
\bottomrule()
\end{longtable}

\begin{quote}
\emph{Bảng 3.2.1 So sánh kiến trúc 3 mô hình thực nghiệm}
\end{quote}

\hypertarget{cux1a1-chux1ebf-chuxfa-uxfd-phuxe2n-cux1ea5p}{%
\subsubsection{3.2.2 Cơ chế chú ý phân
cấp}\label{cux1a1-chux1ebf-chuxfa-uxfd-phuxe2n-cux1ea5p}}

\begin{quote}
Hai tầng chú ý phân cấp. Hệ thống chú ý phân cấp gồm hai tầng hoạt động
bổ sung:

Tầng 1 - Chú ý cục bộ nội khối: SE hoạt động bên trong mỗi khối FR-PDP,
nén thông tin không gian qua GAP thành vector \(1 \times 1 \times C\),
học hàm kích thích qua hai lớp FC để sinh trọng số kênh, khuếch đại kênh
quan trọng và giảm thiểu kênh nhiễu.

Tầng 2 - Chú ý toàn cục ngoại khối: CBAM hoạt động bên ngoài các khối
tại các hooking points, lọc tuần tự theo cả kênh và không gian trên bản
đồ đặc trưng đã tích lũy qua nhiều khối, ở mức trừu tượng cao hơn SE cục
bộ.

Ba cơ sở giữ nguyên SE nội khối:

Cơ sở 1 - Tương thích kiến trúc: Khối FR-PDP dùng chuỗi PW
\(\rightarrow\)DW \(\rightarrow\)PW, trong đó DWConv xử lý mỗi kênh độc
lập trên miền không gian cục bộ \(k \times k\). SE chỉ hoạt động trên
chiều kênh nên bổ sung hoàn toàn cho DWConv. Ngược lại, SAM của CBAM
cũng xử lý miền không gian, gây dư thừa chức năng với DWConv.

Cơ sở 2 - Hiệu quả tham số : SE cực nhẹ (hai lớp FC với tỉ lệ nén
\(r\)). CBAM bổ sung thêm convolution \(7 \times 7\ \)cho SAM. Giữ SE
trong mỗi khối duy trì tính tối giản của mạng siêu nhẹ.

Cơ sở 3 - Phân vai tối ưu: SE là ``bộ lọc kênh nhanh'' cho xử lý nội bộ.
CBAM là ``bộ lọc tổng hợp'' kênh lẫn không gian trên đặc trưng đã qua
nhiều bước xử lý, mang tính ngữ nghĩa cao hơn.

Luồng xử lý đặc trưng. Hình 3.2.2 minh hoa 3 chiến lược chèn cơ chế chú
ý:

\includegraphics[width=2.70833in,height=2.70833in]{vertopal_fe3807bd2608410aa17783721e5463e8/media/image10.jpeg}

\emph{Hình 3.2.2 So sánh 3 mô hình : SE Baseline, CBAM-Local, CBAM-Hook}

Luồng xử lý CBAM-Hook gồm 7 bước:

Bước 1-2: ~Ảnh đầu vào
\(\mathbf{X}_{0} \in \mathbb{R}^{H \times W \times 3}\) qua stem
convolution, sau đó qua các khối FR-PDP pha mở rộng kênh với SE nội khối
và full-residual skip connection:

\(\mathbf{F}_{i} = \text{FR-PDP}_{i}\left( \mathbf{F}_{i - 1} \right) = \mathbf{F}_{i - 1} + \text{SE}\left( \text{PW}_{2}\left( \text{DW}\left( \text{PW}_{1}\left( \mathbf{F}_{i - 1} \right) \right) \right) \right.\ \)
(1)

Bước 3: Tại Hooking Point 1, CBAM toàn cục được áp dụng:

\(\mathbf{F}'_{\text{hook1}} = \text{CBAM}\left( \mathbf{F}_{N} \right) = M_{s}\left( M_{c}\left( \mathbf{F}_{N} \right) \otimes \mathbf{F}_{N} \right) \otimes \left( M_{c}\left( \mathbf{F}_{N} \right) \otimes \mathbf{F}_{N} \right.\ \)
(2)

Trong đó, \(M_{c}\) \(M_{s}\) lần lượt là bản đồ chú ý kênh và không
gian, \(\otimes \ \)là phép nhân theo phần tử.

Bước 4-5: Đặc trưng qua các khối FR-PDP pha thu hẹp, rồi CBAM toàn cục
thứ hai tại Hooking Point 2.

Bước 6: CBAM cuối cùng tại Hooking Point 3 (trước GAP) đảm bảo GAP chỉ
tổng hợp từ vùng không gian mang tính phân biệt cao.

Bước 7: GAP nén bản đồ đặc trưng 3D thành vector 1D, lớp Dense + Softmax
đưa ra phân phối xác suất.

Ưu thế kỳ vọng:

Loại bỏ nhiễu nền vượt trội nhờ ba lớp lọc không gian CBAM dần loại bỏ
vùng nền qua từng pha, hiệu quả trên dữ liệu nền phức tạp.

Khắc phục suy giảm ảnh nhỏ: CBAM-Local làm mất thông tin không gian do
chồng chéo DWConv SAM. CBAM-Hook chỉ lọc không gian ở điểm nối nơi đặc
trưng đã đủ phong phú ngữ nghĩa.

Cân bằng hiệu năng chi phí. Chỉ 3 điểm chèn CBAM thay vì mỗi khối, kiểm
soát tham số bổ sung ở mức tối thiểu.

Các giả thuyết trên được kiểm chứng qua thực nghiệm đối chứng tại mục
3.3.
\end{quote}

\hypertarget{phuxe2n-tuxedch-touxe1n-hux1ecdc-giux1ea3i-quyux1ebft-sux1ef1-chux1ed3ng-chuxe9o-ux111ux1eb7c-trux1b0ng-khuxf4ng-gian-giux1eefa-depthwise-convolution-vuxe0-spatial-attention-module}{%
\subsubsection{3.2.3 Phân tích toán học giải quyết sự chồng chéo đặc
trưng không gian giữa Depthwise Convolution và Spatial Attention
Module}\label{phuxe2n-tuxedch-touxe1n-hux1ecdc-giux1ea3i-quyux1ebft-sux1ef1-chux1ed3ng-chuxe9o-ux111ux1eb7c-trux1b0ng-khuxf4ng-gian-giux1eefa-depthwise-convolution-vuxe0-spatial-attention-module}}

\begin{quote}
Mục này phân tích toán học tại sao CBAM-Local có thể gây suy giảm hiệu
năng, củng cố cơ sở cho chiến lược CBAM-Hook.

Phân tích vùng tiếp nhận
\(\mathbf{F} \in \mathbb{R}^{H \times W \times C}\) ngay sau DWConv
trong khối FR-PDP. Giá trị tại vị trí \((i,j)\)trên kênh \(c\):

\(\mathbf{F}(i,j,c) = \sum_{m = - \left\lfloor k/2 \right\rfloor}^{\left\lfloor k/2 \right\rfloor}{\sum_{n = - \left\lfloor k/2 \right\rfloor}^{\left\lfloor k/2 \right\rfloor}\mathbf{W}_{c}}(m,n) \cdot \mathbf{X}(i + m,j + n,c) + b_{c}\)
(3)

Trong đó, \(\mathbf{W}_{c} \in \mathbb{R}^{k \times k}\ \)là kernel kênh
\(c\ \)( \(k = 3\ \)trong TickNets). DWConv đã mã hóa thông tin không
gian cục bộ \(k \times k\ \)vào \(\mathbf{F}(i,j,c)\).

SAM đặt ngay sau DWConv tạo bản đồ chú ý không gian qua:

Gộp kênh:

\(\mathbf{F}_{\text{avg}}^{s}(i,j) = \frac{1}{C}\sum_{c = 1}^{C}\mathbf{F}(i,j,c),\quad\mathbf{F}_{\text{max}}^{s}(i,j) = \max_{c \in \left\{ 1,\ldots,C \right\}}{\mathbf{F}(i,j,c)}\)
(4)

Tích chập \(7 \times 7\):
\end{quote}

\[M_{s}(i,j) = \sigma\left( \sum_{m = - 3}^{3}{\sum_{n = - 3}^{3}\mathbf{W}^{s}}(m,n) \cdot \left\lbrack \mathbf{F}_{\text{avg}}^{s}(i + m,j + n);\mspace{6mu}\mathbf{F}_{\text{max}}^{s}(i + m,j + n) \right\rbrack + b^{s} \right)\]

\begin{quote}
(5)

Hiện tượng chồng chéo. Thay (3) vào (4):
\end{quote}

\[\mathbf{F}_{\text{avg}}^{s}(i,j) = \frac{1}{C}\sum_{c = 1}^{C}\left\lbrack \sum_{m,n}^{}\mathbf{W}_{c}(m,n) \cdot \mathbf{X}(i + m,j + n,c) + b_{c} \right\rbrack\]

\begin{quote}
(6)

Kết hợp với convolution \(7 \times 7\) của SAM, vùng tiếp cận hiệu dụng
là~:

\(\text{ERF}_{\text{SAM}} = (k - 1) + (7 - 1) + 1 = k + 5\) (7)

Với \(k = 3\): \(\text{ERF}_{\text{SAM}}\)= 8\(\ \)pixel - SAM cục bộ
chỉ học trên thông tin \textbf{cục bộ}, không toàn cục.

Dư thừa nghiêm trọng trên ảnh nhỏ. Với Fashion-MNIST \(28 \times 28\),
sau vài khối FR-PDP có stride \(> 1\), bản đồ đặc trưng
\(H_{l} \times W_{l}\ \)giảm nhanh:

\(\text{Tỉ~lệ~phủ}\) =
\(\frac{\text{ERF}_{\text{SAM}}^{2}}{H_{l} \times W_{l}} \rightarrow 1\quad\text{khi~}H_{l},W_{l} \rightarrow k + 5\)
(8)

Khi tỷ lệ phủ \(\rightarrow 1,\ M_{s}\ \)trở nên đồng nhất - SAM mất khả
năng phân biệt vùng quan trọng và còn đưa thêm nhiễu gradient vào huấn
luyện.

Tại sao CBAM-Hook giải quyết được? Tại các hooking points, mỗi pixel
trên bản đồ đặc trưng đã tích lũy vùng tiếp nhận lớn hơn nhiều nhờ hiệu
ứng chồng tầng:

\(\text{ERF}_{\text{t}\text{í}\text{ch}\text{~}\text{l}\text{ũ}\text{y}}^{(L)} = 1 + \sum_{l = 1}^{L}\left( k_{l} - 1 \right) \cdot \prod_{j = l + 1}^{L}s_{j}\ \)(9)

Với \(L\) đủ lớn,
\(\text{ERF}_{\text{t}\text{í}\text{ch}\text{~}\text{l}\text{ũ}\text{y}}\)
bao phủ phần lớn ảnh đầu vào, giúp SAM phân biệt được ``vùng đối tượng''
và ``vùng nền'' một cách có ý nghĩa.

Kết luận toán học.

CBAM-Local gây dư thừa: SAM sau DWConv bị giới hạn ERF ở mức cục bộ
\(k + 5\ \)pixel, chồng chéo với DWConv. Trên ảnh nhỏ, bản đồ chú ý đồng
nhất và vô nghĩa, gây suy giảm hiệu năng.

CBAM-Hook tối ưu: SAM tại hooking points có ERF tích lũy đủ lớn để hoạt
động trên thông tin ngữ nghĩa toàn cục. SE nội khối vẫn lọc kênh cục bộ
hiệu quả, không xung đột.

Các kết luận này được kiểm chứng bằng thực nghiệm đối chứng tại mục 3.3.
\end{quote}

\begin{enumerate}
\def\labelenumi{\arabic{enumi}.}
\item ~
  \hypertarget{hiux1ec7n-thux1ef1c-giux1ea3i-phuxe1p}{%
  \subsection{HIỆN THỰC GIẢI
  PHÁP:}\label{hiux1ec7n-thux1ef1c-giux1ea3i-phuxe1p}}

  \begin{enumerate}
  \def\labelenumii{\arabic{enumii}.}
  \item ~
    \hypertarget{muxf4i-trux1b0ux1eddng-thux1ef1c-nghiux1ec7m-vuxe0-xuxe2y-dux1ef1ng-muxf4-huxecnh}{%
    \subsubsection{Môi trường thực nghiệm và xây dựng mô
    hình:}\label{muxf4i-trux1b0ux1eddng-thux1ef1c-nghiux1ec7m-vuxe0-xuxe2y-dux1ef1ng-muxf4-huxecnh}}

    \begin{enumerate}
    \def\labelenumiii{\arabic{enumiii}.}
    \item ~
      \hypertarget{muxf4i-trux1b0ux1eddng-thux1ef1c-nghiux1ec7m}{%
      \paragraph{Môi trường thực
      nghiệm}\label{muxf4i-trux1b0ux1eddng-thux1ef1c-nghiux1ec7m}}
    \end{enumerate}
  \end{enumerate}
\end{enumerate}

\begin{quote}
Viết ở đây
\end{quote}

\hypertarget{nhuxf3m-4-bux1ed9-dux1eef-liux1ec7u-thux1ef1c-nghiux1ec7m}{%
\paragraph{Nhóm 4 bộ dữ liệu thực
nghiệm}\label{nhuxf3m-4-bux1ed9-dux1eef-liux1ec7u-thux1ef1c-nghiux1ec7m}}

\begin{quote}
Viết ở đây
\end{quote}

\hypertarget{cuxe1c-phux1b0ux1a1ng-phuxe1p-tiux1ec1n-xux1eed-luxfd-dux1eef-liux1ec7u-vuxe0-tux103ng-cux1b0ux1eddng-dux1eef-liux1ec7u-data-augmentation-cho-tux1eebng-nhuxf3m-ux1ea3nh}{%
\paragraph{Các phương pháp Tiền xử lý dữ liệu và Tăng cường dữ liệu
(Data Augmentation) cho từng nhóm
ảnh}\label{cuxe1c-phux1b0ux1a1ng-phuxe1p-tiux1ec1n-xux1eed-luxfd-dux1eef-liux1ec7u-vuxe0-tux103ng-cux1b0ux1eddng-dux1eef-liux1ec7u-data-augmentation-cho-tux1eebng-nhuxf3m-ux1ea3nh}}

\begin{quote}
Viết ở đây
\end{quote}

\hypertarget{hiux1ec7n-thux1ef1c-huxf3a-3-muxf4-huxecnh-ux111ux1ed1i-chux1ee9ng-thux1ef1c-nghiux1ec7m}{%
\paragraph{Hiện thực hóa 3 mô hình đối chứng thực
nghiệm}\label{hiux1ec7n-thux1ef1c-huxf3a-3-muxf4-huxecnh-ux111ux1ed1i-chux1ee9ng-thux1ef1c-nghiux1ec7m}}

\begin{quote}
Viết ở đây
\end{quote}

\hypertarget{cux1ea5u-huxecnh-quuxe1-truxecnh-huux1ea5n-luyux1ec7n-vuxe0-tux1ed1i-ux1b0u-huxf3a}{%
\paragraph{Cấu hình quá trình huấn luyện và tối ưu
hóa}\label{cux1ea5u-huxecnh-quuxe1-truxecnh-huux1ea5n-luyux1ec7n-vuxe0-tux1ed1i-ux1b0u-huxf3a}}

\begin{quote}
Viết ở đây
\end{quote}

\begin{enumerate}
\def\labelenumi{\arabic{enumi}.}
\setcounter{enumi}{1}
\item ~
  \hypertarget{kux1ebft-quux1ea3-thux1ef1c-nghiux1ec7m-vuxe0-ux111uxe1nh-giuxe1}{%
  \subsubsection{Kết quả thực nghiệm và đánh
  giá:}\label{kux1ebft-quux1ea3-thux1ef1c-nghiux1ec7m-vuxe0-ux111uxe1nh-giuxe1}}

  \begin{enumerate}
  \def\labelenumii{\arabic{enumii}.}
  \item ~
    \hypertarget{cuxe1c-ux111ux1ed9-ux111o-hiux1ec7u-nux103ng-chuux1ea9n-mux1ef1c-accuracy-precision-recall-f1-score-vuxe0-cux1ea5u-truxfac-confusion-matrix}{%
    \paragraph{Các độ đo hiệu năng chuẩn mực Accuracy, Precision,
    Recall, F1-score và cấu trúc Confusion
    Matrix}\label{cuxe1c-ux111ux1ed9-ux111o-hiux1ec7u-nux103ng-chuux1ea9n-mux1ef1c-accuracy-precision-recall-f1-score-vuxe0-cux1ea5u-truxfac-confusion-matrix}}
  \item ~
    \hypertarget{kux1ebft-quux1ea3-thux1ef1c-nghiux1ec7m-vuxe0-biux1ec3u-ux111ux1ed3-lossaccuracy-truxean-tux1eadp-dux1eef-liux1ec7u-chuux1ea9n-fashion-mnist-cifar-10}{%
    \paragraph{Kết quả thực nghiệm và biểu đồ Loss/Accuracy trên tập dữ
    liệu chuẩn Fashion-MNIST,
    CIFAR-10}\label{kux1ebft-quux1ea3-thux1ef1c-nghiux1ec7m-vuxe0-biux1ec3u-ux111ux1ed3-lossaccuracy-truxean-tux1eadp-dux1eef-liux1ec7u-chuux1ea9n-fashion-mnist-cifar-10}}
  \item ~
    \hypertarget{kux1ebft-quux1ea3-thux1ef1c-nghiux1ec7m-vuxe0-biux1ec3u-ux111ux1ed3-lossaccuracy-truxean-tux1eadp-dux1eef-liux1ec7u-thux1ef1c-tiux1ec5n-plantvillage-chest-xray-pneumonia}{%
    \paragraph{Kết quả thực nghiệm và biểu đồ Loss/Accuracy trên tập dữ
    liệu thực tiễn PlantVillage,
    Chest-Xray-Pneumonia}\label{kux1ebft-quux1ea3-thux1ef1c-nghiux1ec7m-vuxe0-biux1ec3u-ux111ux1ed3-lossaccuracy-truxean-tux1eadp-dux1eef-liux1ec7u-thux1ef1c-tiux1ec5n-plantvillage-chest-xray-pneumonia}}
  \item ~
    \hypertarget{phuxe2n-tuxedch-hiux1ec7n-tux1b0ux1ee3ng-suy-giux1ea3m-hiux1ec7u-nux103ng-cux1ee5c-bux1ed9-cbam-local-vuxe0-ux1b0u-thux1ebf-cux1ee7a-giux1ea3i-phuxe1p-ux111ux1ec1-xuux1ea5t-cbam-hook}{%
    \paragraph{Phân tích hiện tượng suy giảm hiệu năng cục bộ CBAM-Local
    và ưu thế của giải pháp đề xuất
    CBAM-Hook}\label{phuxe2n-tuxedch-hiux1ec7n-tux1b0ux1ee3ng-suy-giux1ea3m-hiux1ec7u-nux103ng-cux1ee5c-bux1ed9-cbam-local-vuxe0-ux1b0u-thux1ebf-cux1ee7a-giux1ea3i-phuxe1p-ux111ux1ec1-xuux1ea5t-cbam-hook}}
  \item ~
    \hypertarget{giux1ea3i-thuxedch-muxf4-huxecnh-bux1eb1ng-trux1ef1c-quan-huxf3a-bux1ea3n-ux111ux1ed3-nhiux1ec7t-grad-cam.-sux1ef1-tux1eadp-trung-vuxf9ng-khuxf4ng-gian-ux111ux1eb7c-trux1b0ng-thux1ef1c-tux1ebf}{%
    \paragraph{Giải thích mô hình bằng trực quan hóa bản đồ nhiệt
    Grad-CAM. Sự tập trung vùng không gian đặc trưng thực
    tế}\label{giux1ea3i-thuxedch-muxf4-huxecnh-bux1eb1ng-trux1ef1c-quan-huxf3a-bux1ea3n-ux111ux1ed3-nhiux1ec7t-grad-cam.-sux1ef1-tux1eadp-trung-vuxf9ng-khuxf4ng-gian-ux111ux1eb7c-trux1b0ng-thux1ef1c-tux1ebf}}
  \end{enumerate}
\end{enumerate}

\hypertarget{chux1b0ux1a1ng-4-kux1ebft-quux1ea3-kux1ebft-luux1eadn-vuxe0-kiux1ebfn-nghux1ecb}{%
\section{CHƯƠNG 4: KẾT QUẢ, KẾT LUẬN VÀ KIẾN
NGHỊ}\label{chux1b0ux1a1ng-4-kux1ebft-quux1ea3-kux1ebft-luux1eadn-vuxe0-kiux1ebfn-nghux1ecb}}

\hypertarget{kux1ebft-quux1ea3}{%
\subsection{KẾT QUẢ:}\label{kux1ebft-quux1ea3}}

\begin{quote}
Trong bài tiểu luận này, thông qua các tài liệu về mô hình mạng nơ-ron
tích chập (CNN) và cơ chế chú ý CBAM, học viên đã hiểu kiến trúc, cách
thức hoạt động của mô hình CNN, cùng với đó là cơ chế chú ý và áp dụng
vào mô hình CNN. Học viên đã khảo sát và hoàn thành thực hiện nghiên cứu
việc áp dụng cơ chế chú ý CBAM vào mô hình CNN với pre-trained model
VGG16 trên mặt lý thuyết lẫn thực nghiệm trên hai tập dữ liệu
MNIST-Fashion và Cifar-10. Học viên đã xây dựng được cơ chế chú ý CBAM
áp dụng vào sau một vài lớp của mô hình CNN và thực hiện huấn luyện cùng
với một số mô hình không áp dụng. Từ các số liệu từ hình 3.8,3.9 và bảng
3.3,3.4 so sánh accuracy, precision, recall và f1 score thu được thông
qua thực nghiệm đã cho thấy được mô hình CNN có áp dụng cơ chế chú ý
CBAM đã mang lại kết quả tích cực cho mô hình.
\end{quote}

\hypertarget{kux1ebft-luux1eadn}{%
\subsection{KẾT LUẬN}\label{kux1ebft-luux1eadn}}

\begin{quote}
Qua những thông số thu được từ kết quả thực nghiệm cho thấy kết quả của
mô hình CNN thật sự tốt hơn so với các mô hình máy học truyền thống, đặc
biệt là việc áp dụng cơ chế chú ý CBAM vào mô hình đã cải thiện hiệu
suất cho mô hình thông qua các chỉ số như accuracy, precision, recall và
f1 score. Việc tích hợp CBAM vào mô hình CNN đã chứng minh được hiệu quả
trong việc nâng cao độ chính xác và hiệu suất tổng thể của mô hình phân
loại hình ảnh. Cơ chế Attention giúp mô hình tập trung vào các đặc trưng
và vùng quan trọng, cải thiện khả năng tổng quát hóa. Những kết quả này
khẳng định rằng việc sử dụng cơ chế chú ý à một hướng đi triển vọng
trong việc phát triển các mô hình học sâu cho bài toán phân loại hình
ảnh.
\end{quote}

\hypertarget{kiux1ebfn-nghux1ecb}{%
\subsection{KIẾN NGHỊ}\label{kiux1ebfn-nghux1ecb}}

\begin{quote}
Cơ chế chú ý đã thật sự mang lại hiệu quả tốt cho mô hình CNN, vì vậy
thông qua nghiên cứu này có thể nghiên cứu thêm các cơ chế chú ý khác và
áp dụng nhiều cơ chế chú ý khác nhau để tận dụng các ưu điểm của từng
loại để đạt được hiệu suất cao hơn. Nghiên cứu các phương pháp tối ưu
hóa để giảm thiểu chi phí tính toán của

cơ chế chú ý mà vẫn giữ được hiệu suất cao. Nghiên cứu mở rộng việc áp
dụng cơ chế chú ý vào các nhiệm vụ và các mô hình học sâu.
\end{quote}

\hypertarget{tuxe0i-liux1ec7u-tham-khux1ea3o}{%
\section{TÀI LIỆU THAM KHẢO}\label{tuxe0i-liux1ec7u-tham-khux1ea3o}}

\begin{enumerate}
\def\labelenumi{\arabic{enumi}.}
\item
  Ashish Vaswani, Noam Shazeer, Niki Parmar, Jakob Uszkoreit, Llion
  Jones, Aidan N. Gomez, Łukasz Kaiser, Illia Polosukhin ``Attention is
  all you need'', 2017.
\item
  Trang web \uline{What is Convolutial Neural Network? Begin's Tutorial
  for Machine} \uline{Learning and Deep Learning}, truy cập ngày
  12/03/2024.
\item
  Altenberger, Felix, and Claus Lenz. ``A non-technical survey on deep
  convolutional neural network architectures.'' arXiv preprint arXiv:
  1803.02129 (2018).
\item
  Cao, Changyu et al. ``A convolutional neural network face recognition
  algorithm based on data augmentation.'' (2019)
\item
  Simonyan, K.; Zisserman, A. Very deep convolutional networks for
  large-scale image recognition. arXiv 2014, arXiv:1409.1556..
\item
  Trang web \uline{Thuật toán CNN là gì? Và cấu trúc mạng Convolution
  Neural} \uline{Network}, truy cập ngày 5/4/2024.
\item
  Guo, M.-H., Xu, T.-X., Liu, J.-J., Liu, Z.-N., Jiang, P.-T., Mu,
  T.-J., ... \&
\end{enumerate}

\begin{quote}
Wang, Y. (2022). Attention mechanisms in computer vision: A survey.
Computational Visual Media, 8, 331-368.
\end{quote}

\begin{enumerate}
\def\labelenumi{\arabic{enumi}.}
\setcounter{enumi}{7}
\item
  Trang web \uline{Attention Mechanisms in Computer Vision: CBAM}, truy
  cập ngày 10/05/2024.
\item
  Sanghyun Woo, Jongchan Park, Joon-Young Lee, In So Kweon ``CBAM:
  Convolution Block Attention Module''.
\item
  Jie Hu, Li Shen, Samuel Albanie, Gang Sun, Enhua Wu
  ``Squeeze-and-Excitation Networks''.
\item
  Trang web \uline{CBAM-keras}, truy cập ngày 15/05/2024.
\item
  Trang web \uline{Understanding Transfer Learning for Deep Learning},
  truy cập ngày 15/05/2024.
\item
  Trang web \uline{Fine tuning pre-trained model trong pytorch và áp
  dụng vào} \uline{Visual Saliency Prediction}, truy cập ngày
  16/05/2024.
\end{enumerate}

\end{document}
